\documentclass[a4paper,13pt,oneside]{report}

% =============================================================
% ĐỊNH DẠNG CHUNG -- kế thừa phong cách từ digital_twin_report.tex
% Biên dịch khuyến nghị: pdflatex genmusic_project.tex (2 lần)
% =============================================================
\usepackage{scrextend}
\changefontsizes{13pt}
\usepackage[utf8]{vietnam}
\usepackage[top=2cm,bottom=2cm,left=3.5cm,right=2.5cm]{geometry}
\usepackage{times}
\usepackage{graphicx}
\usepackage{tikz}
\usetikzlibrary{arrows.meta,positioning,calc,fit}
\usepackage{fancybox}
\usepackage{indentfirst}
\usepackage{amsmath,amssymb,amsbsy}
\usepackage{array,tabularx,multirow,booktabs}
\usepackage{enumitem}
\usepackage{titlesec,titletoc}
\usepackage{chngcntr}
\usepackage{float}
\usepackage[font=small,labelfont=bf]{caption}
\usepackage{subcaption}
\usepackage{xcolor}
\usepackage{listings}
\usepackage{setspace}
\usepackage{parskip}
\usepackage{xurl}
\usepackage{ragged2e}
\usepackage{fancyhdr}
\usepackage{hyperref}

\definecolor{codegreen}{rgb}{0,0.5,0}
\definecolor{codegray}{rgb}{0.45,0.45,0.45}
\definecolor{codepurple}{rgb}{0.5,0,0.7}
\definecolor{backcolour}{rgb}{0.96,0.96,0.94}
\lstdefinestyle{mystyle}{
    backgroundcolor=\color{backcolour},
    commentstyle=\color{codegreen},
    keywordstyle=\color{magenta},
    numberstyle=\tiny\color{codegray},
    stringstyle=\color{codepurple},
    basicstyle=\ttfamily\footnotesize,
    breaklines=true,
    captionpos=b,
    keepspaces=true,
    numbers=left,
    numbersep=5pt,
    showstringspaces=false,
    tabsize=2
}
\lstset{style=mystyle}

\def\TITLE{Xây dựng mô hình sinh ca khúc tiếng Việt từ lời bài hát và gợi ý phong cách}
\def\AUTHOR{Đàm Việt Anh, Vũ Thùy Linh, Nguyễn Lê Hùng, Vũ Hồng Quang, Đặng Ngọc Huy}

\hypersetup{
    pdfborder={0 0 0},
    pdftitle={\TITLE},
    pdfauthor={\AUTHOR},
    colorlinks=true,
    linkcolor=black,
    citecolor=black,
    urlcolor=blue
}

\titleformat{\chapter}[hang]{\centering\bfseries}{CHƯƠNG \thechapter.\ }{0pt}{}[]
\titlespacing*{\chapter}{0pt}{-20pt}{20pt}
\titleformat{\section}[hang]{\bfseries}{\thechapter.\arabic{section}\quad}{0pt}{}[]
\titleformat{\subsection}[hang]{\bfseries}{\thechapter.\arabic{section}.\arabic{subsection}\quad}{0pt}{}[]
\titleformat{\subsubsection}[hang]{\bfseries}{\alph{subsubsection})\quad}{0pt}{}[]

\setcounter{secnumdepth}{3}
\counterwithin{figure}{chapter}
\counterwithin{table}{chapter}
\onehalfspacing
\setlength{\parskip}{6pt}
\setlength{\parindent}{15pt}

% Hạn chế dòng mồ côi/góa và việc chỉ một dòng bị đẩy sang trang kế tiếp.
\widowpenalty=10000
\clubpenalty=10000
\displaywidowpenalty=10000
\brokenpenalty=10000
\setlength{\headheight}{18pt}
\emergencystretch=3em
\sloppy

\fancypagestyle{plain}{%
    \fancyhf{}
    \fancyfoot[R]{\thepage}
    \renewcommand{\headrulewidth}{0pt}
    \renewcommand{\footrulewidth}{0pt}
}

\newcommand{\placeholderfigure}[2]{%
    \begin{figure}[H]
        \centering
        \fbox{\begin{minipage}[c][#2][c]{0.88\textwidth}
            \centering\textit{#1}
        \end{minipage}}
    \end{figure}
}

\begin{document}
\hypersetup{pageanchor=false}

% =============================================================
% TRANG BÌA ĐÃ TỐI ƯU -- ĐẠI HỌC BÁCH KHOA HÀ NỘI
% =============================================================
\begin{titlepage}
\thispagestyle{empty}
\begin{center}

{\textbf{\large{ĐẠI HỌC BÁCH KHOA HÀ NỘI}}}\\[0.2cm]
{\textbf{\large{TRƯỜNG CÔNG NGHỆ THÔNG TIN VÀ TRUYỀN THÔNG}}}\\[0.5cm]
% Bạn có thể chèn logo HUST ở đây nếu muốn:
% \includegraphics[width=0.15\textwidth]{logo_hust.png}\\[0.5cm]
\vfill

{\textbf{\huge{BÁO CÁO BÀI TẬP LỚN}}}\\[0.3cm]
{\textbf{\Large{NỀN TẢNG AI TẠO SINH}}}\\[1cm]

{\textbf{\Large{ĐỀ TÀI:}}}\\[0.2cm]
{\textbf{\Large{Xây dựng mô hình sinh ca khúc tiếng Việt\\ từ lời bài hát và gợi ý phong cách}}}\\[1.5cm]

\end{center}

\vfill

% Phần trình bày thông tin Học viên và Giảng viên (chia 2 cột đẹp mắt)
\noindent
\begin{minipage}[t]{0.45\textwidth}
    \textbf{Giảng viên hướng dẫn:}\\
    TS. Dương Quang Huy\\[1cm]
    \textbf{Chương trình đào tạo:}\\
    Thạc sĩ kỹ thuật
\end{minipage}
\hfill
\begin{minipage}[t]{0.5\textwidth}
    \raggedleft
    \textbf{Học viên thực hiện:}\\
    Đàm Việt Anh -- 20252015M\\
    Vũ Thùy Linh -- 20252035M\\
    Nguyễn Lê Hùng -- 20252036M\\
    Vũ Hồng Quang -- 20261114M\\
    Đặng Ngọc Huy -- 20261112M
\end{minipage}

\vfill
\vfill

\begin{center}
{\textbf{\large{HÀ NỘI, 07/2026}}}
\end{center}

\end{titlepage}

\hypersetup{pageanchor=true}
\pagenumbering{roman}

% =============================================================
% TÓM TẮT
% =============================================================
\chapter*{TÓM TẮT}
\addcontentsline{toc}{chapter}{TÓM TẮT}

Đề tài tập trung nghiên cứu và phát triển mô hình trí tuệ nhân tạo sinh ca khúc tiếng Việt tự động từ các tín hiệu điều kiện đầu vào gồm lời ca (văn bản) và phong cách âm nhạc (âm thanh). Hệ thống áp dụng phương pháp chưng cất tri thức (Knowledge Distillation) từ mô hình giáo viên (Teacher) DiffRhythm2 sang mô hình học sinh (Student) MicroDiT nhỏ gọn hơn, sử dụng cơ chế khớp dòng điều kiện (Conditional Flow Matching - CFM) làm cốt lõi để biến đổi nhiễu Gaussian thành biểu diễn phổ/latent nén của giọng hát.

Trong quy trình tiền xử lý tự động, đề tài tích hợp Demucs để phân tách giọng hát và nhạc đệm, Whisper để nhận dạng và căn lề ca từ tiếng Việt kèm mốc thời gian chi tiết, cùng MuQ-MuLan nhằm trích xuất vector phong cách (Audio Style Anchor) 512 chiều. Nhằm khắc phục rào cản ngôn ngữ và đặc trưng thanh điệu phức tạp của tiếng Việt, đề tài tích hợp bộ mã hóa chuyên dụng XPhoneBERT để trích xuất các vector nhúng âm vị (phoneme embeddings). Giải pháp này giúp mô hình nắm bắt chính xác cấu trúc âm tiết, từ đó tối ưu hóa khả năng phát âm rõ chữ và giữ nguyên nhạc tính tự nhiên của ca từ bản địa.

Về mặt hạ tầng, toàn bộ pipeline từ tiền xử lý, huấn luyện đến suy luận đã được kiểm chứng vận hành đồng bộ trên dữ liệu tiếng Việt thật. Về mặt biểu diễn trung gian, đề tài xác nhận rằng đưa mô hình sinh về đúng không gian latent nén (64 chiều, 5~Hz) mà DiffRhythm2 thực sự huấn luyện trên -- thay vì mel-spectrogram trung gian ban đầu -- là hướng đi cho tín hiệu tích cực rõ rệt nhất, sau khi giải quyết được hiện tượng collapse của bộ mã hóa latent tự huấn luyện bằng một bottleneck xác suất (VAE) với lịch trình huấn luyện KL phù hợp. Tuy nhiên, do giới hạn về thời gian và tài nguyên tính toán, cả mô hình sinh trên không gian latent này lẫn quá trình chưng cất tri thức từ teacher đều chưa được huấn luyện tới hội tụ đầy đủ tại thời điểm hoàn thành báo cáo. Kết quả hiện tại cho thấy đây là một hướng tiếp cận khả thi và có tiềm năng, nhưng để đạt chất lượng sinh nhạc ổn định và rõ lời hoàn chỉnh hơn, nút thắt còn lại nhiều khả năng nằm ở quy mô -- dữ liệu huấn luyện, số bước/epoch huấn luyện, hoặc tài nguyên tính toán -- hơn là một giới hạn cấu trúc của phương pháp, và cần được mở rộng thêm trong các nghiên cứu tiếp theo.

\newpage
\renewcommand*\contentsname{MỤC LỤC}
\tableofcontents
\clearpage

\renewcommand{\listfigurename}{DANH MỤC HÌNH VẼ}
\listoffigures
\clearpage

\renewcommand{\listtablename}{DANH MỤC BẢNG BIỂU}
\listoftables
\clearpage

\pagenumbering{arabic}
\pagestyle{fancy}
\fancyhf{}
\fancyhead[R]{\leftmark}
\fancyfoot[R]{\thepage}


% =============================================================
% CHƯƠNG 1: GIỚI THIỆU
% =============================================================
\chapter{GIỚI THIỆU}
\label{chapter:introduction}

\section{Đặt vấn đề}
\label{section:problem}

Sinh ca khúc tự động (Song Generation) là một nhánh nghiên cứu nổi bật của Trí tuệ nhân tạo tạo sinh, trong đó tổng hợp giọng hát (Singing Voice Synthesis - SVS) đóng vai trò then chốt. Để tạo ra một ca khúc tự nhiên, hệ thống phải kiểm soát đồng thời nội dung lời ca, cách phát âm, trường độ, cao độ và nhịp điệu, đồng thời tổ chức chúng hài hòa với phần nhạc đệm.

Đối với tiếng Việt, bài toán này càng phức tạp do đặc trưng của một ngôn ngữ đơn âm tiết và giàu thanh điệu. Cao độ phát âm hàng ngày quyết định ngữ nghĩa của từ, trong khi cao độ giọng hát lại bị chi phối bởi giai điệu bài hát. Do đó, mô hình sinh cần có khả năng dung hòa hai nguồn thông tin này để lời ca vừa bám sát nhạc đệm, vừa giữ được sự rõ ràng về mặt ngữ nghĩa. 

Bên cạnh rào cản ngôn ngữ, việc tổ chức dữ liệu huấn luyện cũng là một thách thức thực tế. Các bản thu âm thông thường là sản phẩm đã hòa âm phối khí trộn lẫn giọng hát và nhạc cụ, đồng thời phần lời ca hiếm khi đi kèm mốc thời gian chi tiết. Quy trình huấn luyện độc lập bắt buộc phải trải qua các bước tiền xử lý phức tạp: phân tách nguồn âm thanh, căn lề ca từ, trích xuất đặc trưng phong cách và chuẩn hóa tín hiệu. Ngoài ra, việc vận hành các mô hình sinh âm nhạc quy mô lớn đòi hỏi tài nguyên tính toán khổng lồ. Để giải quyết các hạn chế này, đề tài hướng tới giải pháp chưng cất tri thức (Knowledge Distillation) nhằm xây dựng một mô hình nhỏ gọn, kế thừa hiệu quả tri thức từ mô hình nền tảng tiền huấn luyện.

\section{Mục tiêu đề tài}
\label{section:objective}

Mục tiêu cốt lõi của đề tài là xây dựng hệ thống có khả năng sinh phổ âm thanh giọng hát chất lượng cao từ các khung điều kiện đầu vào gồm: ca từ tiếng Việt, phần nhạc đệm và đặc trưng phong cách âm nhạc. Hệ thống cần giải quyết tốt sự hòa hợp giữa phát âm thanh điệu bản địa và cấu trúc giai điệu của nhạc đệm.

Đồng thời, đề tài hướng tới tối ưu hóa hiệu năng mô hình sinh thông qua phương pháp chưng cất tri thức. Bằng cách thiết kế cơ chế ánh xạ không gian đặc trưng phù hợp, mô hình học sinh nhỏ gọn có thể tiếp nhận hiệu quả các hướng dẫn từ mô hình giáo viên quy mô lớn, từ đó tăng tốc độ suy luận nhưng vẫn bảo toàn chất lượng âm thanh tự nhiên.

\section{Đóng góp đề tài}
\label{section:scope}

Nghiên cứu tập trung trọng tâm vào thành phần giọng hát tiếng Việt, lấy biểu diễn phổ âm thanh giọng hát làm đầu ra trực tiếp của mô hình sinh. Đóng góp chính của đề tài gồm ba phần: xây dựng một quy trình tự động tiền xử lý dữ liệu âm thanh thô; triển khai và huấn luyện mạng nơ-ron dự đoán biểu diễn phổ giọng hát dựa trên kỹ thuật chưng cất tri thức; và đánh giá thực nghiệm hiệu năng hệ thống qua các chỉ số kỹ thuật cùng chất lượng cảm tính.

\section{Bố cục báo cáo}
\label{section:report-structure}

Nội dung báo cáo được cấu trúc thành năm chương. Chương 2 trình bày tổng quan các nghiên cứu tiền đề và nền tảng lý thuyết cốt lõi về mô hình sinh âm thanh, cơ chế khớp dòng chảy, kiến trúc mạng nơ-ron, phương pháp chưng cất tri thức và bộ giải mã âm thanh. Chương 3 mô tả chi tiết kiến trúc tổng quan của hệ thống, quy trình tiền xử lý, thiết kế mạng sinh giọng hát, cơ chế huấn luyện và giải pháp chưng cất đặc trưng thích ứng. Chương 4 trình bày thiết lập thực nghiệm, nguồn dữ liệu, các chỉ số đánh giá, kết quả của từng thí nghiệm và thảo luận theo câu hỏi nghiên cứu. Chương 5 tổng kết các kết quả đạt được và đề xuất định hướng nghiên cứu tiếp theo.


% =============================================================
% CHƯƠNG 2
% =============================================================
\chapter{CƠ SỞ LÝ THUYẾT}
\label{chapter:related}

Sinh âm nhạc từ văn bản là bài toán ánh xạ một mô tả ngôn ngữ tự nhiên sang tín hiệu âm thanh có cấu trúc. Khác với sinh một âm thanh sự kiện ngắn, một bài hát phải duy trì đồng thời nội dung ca từ, giai điệu, nhịp, hòa âm, màu sắc nhạc cụ và cấu trúc theo thời gian. Mô hình vì vậy không chỉ tạo tín hiệu hợp lệ, mà còn phải bảo đảm quan hệ giữa câu hát và vùng thời gian tương ứng, đồng thời giữ tính nhất quán trên chuỗi dài.

Các hệ thống hiện đại thường không học trực tiếp trên từng mẫu waveform. Một đoạn 10 giây ở 24~kHz đã chứa 240.000 mẫu, gây chi phí rất lớn cho transformer. Âm thanh do đó được đổi sang biểu diễn nén hơn như token neural codec, latent hoặc mel-spectrogram; mô hình sinh học phân phối của biểu diễn trung gian, sau đó vocoder khôi phục waveform. MusicGen đại diện cho hướng tự hồi quy trên mã codec rời rạc \cite{musicgen}. AudioLDM sử dụng latent diffusion cùng biểu diễn âm thanh - văn bản tương phản \cite{audioldm}. DiffRhythm hướng tới bài hát có lời, kết hợp lyric, phong cách và latent diffusion để sinh đồng thời giọng hát cùng phần đệm \cite{diffrhythm}. Đề tài kế thừa hướng liên tục của DiffRhythm nhưng xây dựng một student nhỏ, huấn luyện bằng Conditional Flow Matching (CFM) và tiếp nhận tri thức từ DiffRhythm2.


\section{Biểu diễn mel-spectrogram của tín hiệu âm thanh}

Mel-spectrogram là biểu diễn tần số-thời gian phổ biến trong xử lý âm thanh, giúp nén sóng âm (waveform) thành ma trận hai chiều mà vẫn bảo toàn cấu trúc phổ biến thiên theo thời gian. 

Quá trình trích xuất bắt đầu bằng việc chia tín hiệu rời rạc $x[n]$ thành các khung ngắn và áp dụng phép biến biến đổi Fourier thời gian ngắn (STFT) tại khung $t$, tần số $k$:
\begin{equation}
 X(t,k)=\sum_{n=0}^{N-1} x[n+tH]w[n]\exp\left(-j\frac{2\pi kn}{N}\right),
\end{equation}
với $N$ là kích thước FFT, $H$ là hop length và $w[n]$ là cửa sổ phân tích. Phổ công suất $|X(t,k)|^2$ sau đó được chiếu qua ngân hàng bộ lọc tam giác trên thang đo Mel. 

Thang đo Mel là một mô hình tâm sinh vật lý đã được kiểm chứng thực nghiệm, phản ánh chính xác cơ chế cảm nhận cao độ phi tuyến tính của tai người \cite{mel}. Nghiên cứu sinh lý học chỉ ra rằng hệ thống thính giác của con người nhạy bén hơn với các thay đổi ở dải tần số thấp và kém nhạy dần ở dải tần số cao. Do đó, việc áp dụng ngân hàng bộ lọc Mel giúp phân bổ độ phân giải tần số một cách tối ưu - dày đặc ở dải âm trầm và thưa dần ở dải âm cao - đồng thời làm nổi bật các đặc trưng âm sắc tiệm cận với nhận thức thực tế của con người \cite{mfcc}. Phổ log-mel tương ứng được tính bởi:
\begin{equation}
 M(t,r)=\log\left(\epsilon+\sum_k A_{r,k}|X(t,k)|^2\right),
\end{equation}
trong đó $A_{r,k}$ là trọng số bộ lọc Mel thứ $r$ tại tần số $k$, và $\epsilon$ là hằng số nhỏ để tránh lỗi logarit của $0$.

Mặc dù là đặc trưng tối ưu cho học sâu nhờ tính tương thích sinh học này, việc khôi phục trực tiếp dạng sóng từ mel-spectrogram là một bài toán xấp xỉ ước lượng do sự hao hụt thông tin không thể đảo ngược: (1) loại bỏ hoàn toàn thông tin pha và (2) phép chiếu Mel làm gộp nhiều tần số tuyến tính ở dải cao. 

Để giải quyết, Vocos \cite{vocos} được lựa chọn làm neural vocoder thế hệ mới dựa trên mạng GAN. Khác với các kiến trúc truyền thống (như HiFi-GAN hay BigVGAN) sử dụng chuỗi lớp tích chập chuyển vị để nâng bậc mẫu tuần tự, Vocos giữ nguyên độ phân giải thời gian xuyên suốt (isotropic) và tái tạo dạng sóng hiệu quả thông qua phép biến đổi Fourier ngược (iSTFT).

\section{Biểu diễn lời ca bằng XPhoneBERT}
\label{sec:bieu-dien-loi-ca-xphonebert}

Để tổng hợp giọng hát tự nhiên, hệ thống cần biểu diễn chính xác thông tin ngữ âm từ ca từ. Trong ngôn ngữ học, chữ viết thô được cấu thành từ các \textbf{âm tự} (grapheme - đơn vị chữ viết nhỏ nhất như ``a'', ``b'', hoặc tổ hợp ``ngh''), trong khi cách phát âm thực tế lại được quyết định bởi các \textbf{âm vị} (phoneme - đơn vị âm thanh nhỏ nhất giúp phân biệt nghĩa). 

Đối với tiếng Việt, việc sử dụng trực tiếp chuỗi âm tự bộc lộ nhiều hạn chế do quan hệ bất đối xứng giữa chữ viết và phát âm thực tế. Bảng \ref{tab:vietnamese_consonants} minh họa sự không đồng nhất này thông qua hệ thống phụ âm tiếng Việt, nơi nhiều âm tự khác nhau có thể biểu diễn chung một âm vị ngữ âm học (ví dụ: các âm tự ``c'', ``k'', ``q'' đều phát âm chung âm vị $/k/$, hoặc ``ng'', ``ngh'' đều phát âm chung âm vị $/\eta/$).

Ngoài ra, tiếng Việt là ngôn ngữ giàu thanh điệu; các dấu thanh đóng vai trò là các âm vị siêu đoạn tính điều khiển trực tiếp tần số cơ bản ($F_0$) và đường cong cao độ, nhưng lại chỉ được biểu thị tối giản trên văn bản viết. Do đó, để loại bỏ sự phụ thuộc vào các biến thể chính tả và đưa không gian biểu diễn về gần hơn với miền âm học, ca từ thô $y$ bắt buộc phải được chuyển đổi thành chuỗi âm vị thông qua bộ chuyển đổi tự vị -- âm vị (G2P):
\begin{equation}
 p = \operatorname{G2P}(y) = (p_1, p_2, \ldots, p_L),
 \label{eq:lyric_g2p}
\end{equation}
trong đó $p_i$ đại diện cho một đơn vị âm vị và $L$ là độ dài chuỗi kết quả.

\begin{table}[H]
\centering
\caption{Ví dụ minh họa sự bất đối xứng giữa âm vị và âm tự trong hệ thống phụ âm tiếng Việt}
\label{tab:vietnamese_consonants}
\small
\renewcommand{\arraystretch}{0.92}
\setlength{\tabcolsep}{2.5pt}
\begin{tabular}{|c|c|c|l|}
\hline
\textbf{STT} & \textbf{Âm vị (Ký hiệu phiên âm)} & \textbf{Âm tự (Chữ viết thể hiện)} & \textbf{Cách phát âm minh họa} \\ \hline
1  & $/b/$ & b & bờ \\ \hline
2  & $/k/$ & c, k hoặc q & cờ, ca hoặc quờ \\ \hline
3  & $/d/$ & đ & đờ \\ \hline
4  & $/y/$ & g hoặc gh & gờ hoặc ghờ \\ \hline
5  & $/h/$ & h & hờ \\ \hline
6  & $/l/$ & l & lờ \\ \hline
7  & $/m/$ & m & mờ \\ \hline
8  & $/n/$ & n & nờ \\ \hline
9  & $/p/$ & p & pờ \\ \hline
10 & $/z/$ & r & rờ \\ \hline
11 & $/s/$ & s & sờ \\ \hline
12 & $/t/$ & t & tờ \\ \hline
13 & $/v/$ & v & vờ \\ \hline
14 & $/s/$ & x & xờ \\ \hline
15 & $/c/$ hoặc $/t/$ & ch & chờ \\ \hline
16 & $/z/$ & d hoặc gi & dờ \\ \hline
17 & $/x/$ hoặc $/k^h/$ & kh & khờ \\ \hline
18 & $/n/$ & ng hoặc ngh & ngờ hoặc nghờ \\ \hline
19 & $/n/$ & nh & nhờ \\ \hline
20 & $/f/$ & ph & phờ \\ \hline
21 & $/t'/ $ hoặc $/t^h/$ & th & thờ \\ \hline
22 & $/t/$ hoặc $/t\int/$ & tr & trò \\ \hline
\end{tabular}
\end{table}



Tuy nhiên, các âm vị đơn lẻ trong chuỗi mới chỉ là các đặc trưng tĩnh, chưa thể hiện được sự biến thiên liên tục và ảnh hưởng lẫn nhau của các âm tiết đứng cạnh nhau khi hát. Để giải quyết vấn đề này, đề tài sử dụng mô hình ngôn ngữ tiền huấn luyện quy mô lớn XPhoneBERT \cite{xphonebert}. Là một mô hình Masked Language Model hoạt động ở cấp độ âm vị và được tối ưu hóa cho các ngôn ngữ đơn âm tiết, đa thanh điệu như tiếng Việt, mô hình mã hóa đầu vào thành trạng thái ẩn, tích hợp chặt chẽ đặc trưng phát âm ngữ cảnh của toàn bộ câu hát, cung cấp ràng buộc về phát âm để mô hình sinh ở các giai đoạn sau thực hiện căn chỉnh thời gian và tái tạo tiếng hát tròn vành rõ chữ.

\section{Biểu diễn phong cách âm nhạc với MuQ-MuLan}

Mô hình MuLan \cite{mulan} giải quyết khoảng cách ngữ nghĩa giữa mô tả văn bản (như thể loại, nhịp điệu, tốc độ) và tín hiệu âm thanh thông qua việc huấn luyện tương phản song song hai bộ mã hóa (audio/text encoder) trên cùng một không gian biểu diễn chung. Kế thừa nguyên lý đó, MuQ-MuLan \cite{muq} kết hợp bộ mã hóa âm nhạc tự giám sát MuQ với cơ chế học tương phản đa phương thức để trích xuất vector phong cách toàn cục (style embedding). Trong đề tài này, vector biểu diễn đóng vai trò là "neo phong cách", tóm tắt các đặc trưng phi ngôn ngữ cốt lõi như màu sắc âm thanh, cấu trúc phối khí và nhịp điệu từ bản nhạc tham chiếu.

Để tối ưu hóa quá trình sinh giọng hát, hệ thống phân tách thông tin điều kiện thành hai nhánh độc lập: ngôn ngữ (cục bộ) và phong cách (toàn cục). Bộ mã hóa văn bản chịu trách nhiệm kiểm soát nội dung âm vị học chi tiết theo trục thời gian, trong khi MuQ-MuLan cung cấp định hướng thẩm mỹ âm nhạc tổng thể. Cấu trúc hai nhánh song hành này giải quyết triệt để giới hạn của các phương pháp đơn lẻ: nhánh ngôn ngữ đảm bảo lời ca được phát âm rõ ràng, chính xác; còn nhánh phong cách đảm bảo giọng hát được đặt trong một không gian hòa âm phối khí đồng nhất và tự nhiên.

\section{Mô hình khuếch tán DDPM và Conditional Flow Matching}
\label{sec:ddpm_cfm}

Mô hình khuếch tán và mô hình khớp dòng đều hướng tới việc biến đổi một phân phối nguồn đơn giản (thường là Gaussian) thành phân phối dữ liệu đích, nhưng khác nhau về cách tham số hóa quá trình biến đổi. Trong khi DDPM mô hình hóa quá trình khử nhiễu ngược qua một chuỗi Markov rời rạc \cite{ddpm}, Flow Matching hồi quy trực tiếp trường vector của một dòng chuẩn hóa liên tục (CNF) \cite{flowmatching}. Rectified Flow có liên hệ mật thiết với Flow Matching nhưng nhấn mạnh việc làm thẳng quỹ đạo vận chuyển giữa hai đầu mút \cite{rectifiedflow}.

Trong DDPM, quá trình khuếch tán thuận biến dữ liệu sạch $Z$ thành nhiễu qua $K$ bước. Trạng thái tại bước $k$ bất kỳ được lấy mẫu trực tiếp qua công thức tái tham số hóa \cite{ddpm}:
\begin{equation}
 Z_k = \sqrt{\bar{\alpha}_k} Z + \sqrt{1-\bar{\alpha}_k} \epsilon, \quad \epsilon \sim \mathcal{N}(0, \mathbf{I}),
 \label{eq:ddpm_reparam}
\end{equation}
trong đó $\bar{\alpha}_k$ là hệ số nhiễu tích lũy. Mạng nơ-ron $\epsilon_\theta$ được huấn luyện để dự đoán nhiễu $\epsilon$ từ trạng thái nhiễu $Z_k$ và điều kiện $c$ bằng cách tối ưu hóa hàm mất mát:
\begin{equation}
 \mathcal{L}_{\mathrm{DDPM}} = \mathbb{E}_{Z,\epsilon,k} \left[ \left\| \epsilon - \epsilon_\theta(Z_k, k, c) \right\|_2^2 \right].
 \label{eq:ddpm_loss}
\end{equation}
Hạn chế lớn nhất của DDPM là quá trình suy luận mang tính tuần tự qua $K$ bước rời rạc, gây tốn kém chi phí tính toán khi số bước lấy mẫu lớn \cite{ddpm}.

Để khắc phục nhược điểm này, Flow Matching định nghĩa một dòng liên tục $\{p_t\}_{t\in[0,1]}$ được sinh bởi phương trình vi phân thường (ODE):
\begin{equation}
 \frac{\mathrm{d}Z_t}{\mathrm{d}t} = u_t(Z_t,c), \quad Z_0 \sim p_0, \quad Z_1 \sim p_{\mathrm{data}}(\cdot \mid c).
 \label{eq:fm_ode}
\end{equation}
Do việc tính toán trực tiếp trường vector biên $u_t$ là bất khả thi, Conditional Flow Matching (CFM) sử dụng đường nội suy thẳng nối giữa hai đầu mút $Z_0 \sim \mathcal{N}(0, \mathbf{I})$ và $Z_1 \sim p_{\mathrm{data}}$ \cite{flowmatching}. Tại thời điểm $t \sim \mathcal{U}(0,1)$, trạng thái trung gian và vận tốc có điều kiện tương ứng được xác định đơn giản bởi:
\begin{equation}
 Z_t = (1-t)Z_0 + tZ_1 \quad \Rightarrow \quad \frac{\mathrm{d}Z_t}{\mathrm{d}t} = Z_1 - Z_0.
 \label{eq:cfm_interpolation}
\end{equation}
Mạng $v_\theta$ (nhận điều kiện $c$ gồm lời ca, phong cách và nhạc đệm) được huấn luyện để dự đoán trường vận tốc này thông qua hàm mất mát CFM \cite{flowmatching}:
\begin{equation}
 \mathcal{L}_{\mathrm{CFM}} = \mathbb{E}_{c,Z_1,Z_0,t} \left[ \left\| v_\theta(Z_t, t, c) - (Z_1 - Z_0) \right\|_2^2 \right].
 \label{eq:cfm_loss}
\end{equation}
Nghiệm tối ưu của mạng chính là kỳ vọng có điều kiện $v^*(z,t,c) = \mathbb{E}[Z_1 - Z_0 \mid Z_t = z, t, c]$. Khi suy luận, ta chỉ cần giải ODE dọc theo trường vận tốc đã học từ $t=0$ đến $t=1$ bằng các bộ giải như Euler \cite{flowmatching}. Nhờ quỹ đạo vận chuyển gần như thẳng, CFM cho phép giảm tối đa số bước tích phân mà vẫn đảm bảo chất lượng mẫu vượt trội so với các mô hình khuếch tán truyền thống.

\section{Mô hình sinh nhạc DiffRhythm 2}
\label{sec:diffrhythm2_framework}

DiffRhythm 2 là hệ thống end-to-end sinh bài hát đầy đủ cấu trúc (gồm cả giọng hát và nhạc đệm) có độ dài lên tới 210 giây từ ca từ và mô tả phong cách \cite{diffrhythm2}. Hệ thống áp dụng kiến trúc bán tự hồi quy dựa trên cơ chế khớp dòng khối (\textbf{Block Flow Matching - BFM}) \cite{diffrhythm2}. Chuỗi biểu diễn ẩn (latent) được chia thành các khối có độ dài cố định $b=10$ (tương đương 2 giây âm thanh ở tần số latent 5~Hz). Trong nội bộ từng khối, các vị trí được sinh song song không tự hồi quy (NAR) bằng Flow Matching; giữa các khối, mối quan hệ nhân quả được bảo toàn qua attention mask, kết hợp tối ưu khả năng xử lý song song với việc duy trì ngữ cảnh dài hạn \cite{flowmatching, diffrhythm2}.

Để giảm chiều dài chuỗi đầu vào cho Diffusion Transformer, DiffRhythm 2 sử dụng một bộ Music VAE hoạt động ở tần số latent cực thấp (5~Hz) \cite{diffrhythm2}. Bộ mã hóa nhận âm thanh đầu vào 24~kHz và bộ giải mã tái tạo âm thanh ở tần số 48~kHz, đạt tỷ lệ nén lần lượt là 4.800 và 9.600 lần \cite{diffrhythm2}. Bộ giải mã tích hợp một khối Transformer bổ trợ cùng backbone BigVGAN \cite{bigvgan}. Toàn bộ Music VAE được tối ưu hóa bằng tổ hợp mel loss đa thang đo và STFT loss đa độ phân giải nhằm giữ nguyên vẹn thông tin của cả giọng hát lẫn nhạc đệm \cite{diffrhythm2}.

Mô hình sinh chính của hệ thống là một Diffusion Transformer xấp xỉ 1 tỷ tham số \cite{dit, diffrhythm2}. Khi dự đoán khối thứ $i$, đầu vào bao gồm latent nhiễu của khối hiện tại $z_t^i$ ($t \sim \mathcal{U}(0,1)$), các khối sạch lịch sử $z^{<i}$ (được gán timestep cố định $1.0$ để phân biệt), cùng chuỗi ca từ và đặc trưng phong cách từ MuQ-MuLan \cite{muq, diffrhythm2}. Công trình cũng giới thiệu hàm mất mát **Stochastic Block REPA** nhằm căn chỉnh biểu diễn ẩn của mô hình sinh với đặc trưng tự giám sát (SSL) toàn cục từ một tập con các khối được lấy mẫu ngẫu nhiên \cite{diffrhythm2}. Sau giai đoạn tiền huấn luyện, mô hình được tinh chỉnh bằng kỹ thuật tối ưu hóa ưu tiên cặp chéo (\textbf{Cross-Pair Preference Optimization - CPPO}) để cân bằng đa mục tiêu giữa tính nhạc, độ khớp lời và chất lượng âm thanh \cite{dpo, diffrhythm2}. Khi suy luận, cấu trúc khối cho phép áp dụng KV cache và FlexAttention, đạt hệ số thời gian thực (RTF) khoảng 0.213 trên cấu hình thử nghiệm của tác giả \cite{diffrhythm2}.

Trong khuôn khổ đề tài này, DiffRhythm 2 đóng vai trò là một mô hình Teacher đã được đóng băng tham số. Một mô hình con sẽ được học tín hiệu chuyển giao là trường vận tốc (velocity), bắt chước hướng di chuyển tương tự như trường vector của mô hình Teacher.


% =============================================================
% CHƯƠNG 3
% =============================================================
\chapter{PHƯƠNG PHÁP}
\label{chapter:method}


\section{Chuẩn bị dữ liệu}
\label{sec:data_preparation}

Dữ liệu có giám sát đầy đủ cho bài toán sinh giọng hát tiếng Việt còn hạn chế, trong khi nguồn dữ liệu thực tế thường chỉ là bản phối hoàn chỉnh, không có giọng hát tách riêng và không đi kèm căn chỉnh ca từ. Do đó, đề tài xây dựng quy trình chuẩn bị dữ liệu theo hướng giám sát yếu, phù hợp với định dạng dữ liệu của DiffRhythm2 \cite{diffrhythm2}. Từ mỗi bài hát thô, quy trình tạo ra một mẫu gồm mel-spectrogram của giọng hát mục tiêu, mel-spectrogram của phần đệm, ca từ có mốc thời gian và vector phong cách toàn cục. Các nhãn thu được từ mô hình tiền huấn luyện chỉ mang tính xấp xỉ và phụ thuộc vào chất lượng của các mô hình sử dụng, do đó việc hậu kiểm thủ công cần được thực hiện trong tương lai.

Đối với thành phần tách nguồn, mỗi bản phối đầu vào được chuẩn hóa về tần số lấy mẫu 24~kHz và được đưa qua HTDemucs \cite{htdemucs}. Mô hình sẽ tách âm thanh ban đầu thành hai thành phần, bao gồm giọng hát và nhạc đệm. Giọng hát tách được đóng vai trò mục tiêu học của mô hình sinh, trong khi phần đệm được giữ làm điều kiện âm nhạc theo thời gian. Điều kiện này mang thông tin về nhịp, hòa âm, chuyển đoạn và khoảng nghỉ, nhờ đó cung cấp ngữ cảnh cục bộ mà một mô tả phong cách ngắn không thể biểu diễn đầy đủ.

Đối với thành phần lời, Whisper \cite{whisper} được áp dụng lên phần giọng hát để nhận dạng tiếng Việt và sinh các đoạn văn bản có mốc thời gian . Mỗi đoạn gồm nội dung lời và cặp thời điểm bắt đầu--kết thúc trong tín hiệu âm thanh. Khi tạo một đoạn huấn luyện, chỉ các đoạn lời giao với cửa sổ thời gian tương ứng được sử dụng làm điều kiện văn bản; các đoạn này được ghép theo thứ tự thời gian trước khi đưa vào bộ mã hóa ngôn ngữ. Cách tổ chức này giúp điều kiện lời bám với đoạn âm thanh đang học, đồng thời tránh việc mô hình phải liên hệ toàn bộ lời bài hát với một đoạn giọng hát ngắn. Nếu cửa sổ âm thanh không chứa lời, điều kiện văn bản được xem là rỗng và mô hình vẫn học từ phần đệm cùng vector phong cách.

Đối với thành phần phổ âm thanh, cả giọng hát mục tiêu và phần đệm đều được biến đổi sang mel-spectrogram trong cùng miền đặc trưng với Vocos \cite{vocos}. Cấu hình sử dụng tần số lấy mẫu 24~kHz, 100 mel bins, kích thước FFT bằng 1024 mẫu và hop length bằng 256 mẫu. Mỗi đoạn huấn luyện gồm 128 khung mel, tương ứng thời lượng xấp xỉ $128\times256/24000\approx1{,}36$~giây. Như vậy, mel mục tiêu $x_1$ và mel điều kiện $b$ có cùng kích thước $100\times128$ cho mỗi mẫu. Việc thống nhất cấu hình mel ngay từ tiền xử lý giúp đầu ra của mô hình có thể được giải mã trực tiếp bằng Vocos mà không cần nội suy lại trục mel.

Đối với thành phần phong cách, bản phối gốc được mã hóa bằng MuQ-MuLan thành
vector 512 chiều \cite{muq}. Vector này được trích từ đoạn âm thanh tham chiếu
24~kHz, ưu tiên tối đa 10 giây đầu của bài hát, để tóm tắt đặc trưng chung của
bản nhạc như màu âm, phối khí, năng lượng và xu hướng thể loại. Trong mẫu huấn
luyện, vector phong cách $s\in\mathbb{R}^{512}$ không mô tả chi tiết từng thời
điểm, mà đóng vai trò như ngữ cảnh toàn cục của bài hát. Ngược lại, mel của phần
đệm cung cấp diễn tiến cục bộ theo thời gian. Hai nguồn điều kiện này bổ sung cho
nhau: phần đệm cho biết bài hát đang diễn ra như thế nào, còn vector phong cách
cho biết bài hát có màu sắc âm nhạc tổng quát ra sao.

Sau tiền xử lý, mỗi đoạn huấn luyện được biểu diễn bởi bộ $(x_1,b,l,m,s)$.
Trong đó, $x_1\in\mathbb{R}^{100\times128}$ là mel-spectrogram mục tiêu của
giọng hát; $b\in\mathbb{R}^{100\times128}$ là mel-spectrogram của phần đệm trong
cùng khoảng thời gian; $l$ là chuỗi mã hóa ca từ; $m$ là mặt nạ đánh dấu các vị
trí token hợp lệ; và $s\in\mathbb{R}^{512}$ là vector phong cách toàn cục. Cách
biểu diễn này giúp mỗi mẫu giữ đồng thời ba loại thông tin cần thiết cho mô hình:
nội dung lời, ngữ cảnh âm nhạc theo thời gian và phong cách tổng quát của bài hát.

\section{Huấn luyện CFM}
\label{sec:microdit_architecture}

\subsection{Phát biểu bài toán và ký hiệu}

Dữ liệu huấn luyện được biểu diễn bởi tập
\begin{equation}
 \mathcal{D}=\{(y_i,a_i^{\mathrm{voc}},a_i^{\mathrm{acc}})\}_{i=1}^{N},
\end{equation}
trong đó $y_i$ là ca từ có mốc thời gian, $a_i^{\mathrm{voc}}$ là phần giọng hát và $a_i^{\mathrm{acc}}$ là phần nhạc đệm. Sau tiền xử lý, mỗi mẫu gồm mel vocal $x_1^{\mathrm{voc}}\in\mathbb{R}^{F\times T}$, mel backing $x_1^{\mathrm{acc}}\in\mathbb{R}^{F\times T}$, chuỗi ca từ $l$ và style embedding $s\in\mathbb{R}^{512}$. Mục tiêu huấn luyện thực tế $x_1$ là **hỗn hợp cả bài** (giọng hát cộng nhạc đệm), tái tạo bằng cách cộng năng lượng biên độ tuyến tính của hai mel rồi lấy log lại:
\begin{equation}
 x_1=\log\!\left(\exp(x_1^{\mathrm{voc}})+\exp(x_1^{\mathrm{acc}})\right)
 \label{eq:reconstruct_full_mix}
\end{equation}
(áp dụng trên giá trị mel đã denormalize). Hệ thống cần học trường velocity:
\begin{equation}
 v_\theta:\;(x_t,t,l,s)\mapsto \hat{u}_t\in\mathbb{R}^{F\times T},
\end{equation}
để biến nhiễu Gaussian thành mel của **hỗn hợp cả bài** phù hợp điều kiện lời và phong cách. Trong cấu hình student, $F=100$ và mỗi đoạn huấn luyện có $T=128$ khung mel (mặc định hiện tại là $384$), tương ứng khoảng 1,36--4,1 giây ở 24~kHz với hop length 256.

\textbf{Chú ý quan trọng về vai trò của backing}: $x_1^{\mathrm{acc}}$ (backing) chỉ được dùng để \emph{xây dựng mục tiêu huấn luyện} $x_1$ qua công thức \eqref{eq:reconstruct_full_mix} — nó **không** còn là một tensor điều kiện đưa vào mạng $v_\theta$ (một số phiên bản kiến trúc trước đây của dự án có nhánh này, đã bị loại bỏ, xem \S\ref{sec:microdit_architecture} tiếp theo). Do đó model phải học sinh ra cả bài (giọng hát và nhạc đệm) từ trạng thái không điều kiện theo backing, chỉ dựa vào ca từ và style embedding toàn cục. Để tránh model "lãng quên" tín hiệu giọng hát (yếu hơn, thưa hơn nhạc đệm) khi tối ưu trực tiếp trên hỗn hợp, một đầu ra phụ trợ dự đoán riêng velocity của vocal (\emph{vocal auxiliary head}) được huấn luyện đồng thời — chi tiết ở \S\ref{sec:microdit_architecture}.

\subsection{Kiến trúc tổng thể của hệ thống}

Hình~\ref{fig:pipeline} mô tả kiến trúc tổng thể của hệ thống. Kiến trúc không phải một chuỗi xử lý tuyến tính đơn giản mà gồm
ba lớp chức năng. Lớp dữ liệu tạo ra các điều kiện đồng bộ từ bài hát thô; lớp mô hình
học trường vận tốc bằng CFM hoặc bằng chưng cất; lớp suy luận tích phân trường vận tốc
và giải mã mel thành waveform. Giữa các lớp là những hợp đồng dữ liệu có thể kiểm
tra độc lập.

\begin{figure}[H]
\centering
\begin{tikzpicture}[
    >=Latex,
    source/.style={draw, rounded corners=2pt, fill=gray!10,
        minimum height=0.8cm, text width=2.8cm, align=center, font=\scriptsize},
    process/.style={draw, rounded corners=2pt, fill=blue!7,
        minimum height=0.85cm, text width=3.2cm, align=center, font=\scriptsize},
    artifact/.style={draw, rounded corners=2pt, fill=green!8,
        minimum height=0.85cm, text width=4.1cm, align=center, font=\scriptsize},
    model/.style={draw, rounded corners=2pt, fill=orange!12,
        minimum height=0.9cm, text width=3.5cm, align=center, font=\scriptsize},
    optional/.style={draw, dashed, rounded corners=2pt, fill=purple!5,
        minimum height=0.85cm, text width=2.8cm, align=center, font=\scriptsize},
    flow/.style={-{Latex[length=2mm]}, line width=0.55pt},
    optflow/.style={-{Latex[length=2mm]}, dashed, line width=0.5pt}
]
\node[source] (lyric) at (-4.8,0) {Ca từ và mô tả phong cách};
\node[source] (audio) at (0,0) {Tập tin WAV/MP3};
\node[process] (norm) at (-4.8,-1.35) {Chuẩn hóa văn bản và mốc thời gian};
\node[process] (demucs) at (-0.8,-1.35) {Demucs: tách vocal và backing};
\node[process] (muq) at (3.4,-1.35) {MuQ-MuLan: style embedding 512 chiều};
\node[process] (whisper) at (-2.6,-2.8) {Whisper: transcript và segment timestamps};
\node[process] (mel) at (1.5,-2.8) {Vocos feature extractor: mel 100 bins};
\node[artifact] (records) at (-0.5,-4.25)
    {Bản ghi hợp nhất + vocal/backing mel + style embedding};
\node[artifact] (validation) at (-0.5,-5.55)
    {Kiểm tra tensor, metadata và lấy crop đồng bộ};
\node[model] (self) at (-3.0,-7.0)
    {MicroDiT + CFM\\(tự huấn luyện)};
\node[model] (distill) at (2.0,-7.0)
    {DiffRhythm2 teacher $\rightarrow$ MicroDiT\\(chưng cất)};
\node[optional] (kaggle) at (5.2,-5.55)
    {Kaggle backend tùy chọn\\GPU train/inference};
\node[process] (sample) at (-0.5,-8.45)
    {Euler ODE sampling + Vocos 24 kHz};
\node[artifact] (output) at (-0.5,-9.75)
    {Waveform WAV/MP3 và báo cáo đánh giá};

\draw[flow] (lyric) -- (norm);
\draw[flow] (audio) -- (demucs);
\draw[flow] (audio) -| (muq);
\draw[flow] (demucs) -- (whisper);
\draw[flow] (demucs) -- (mel);
\draw[flow] (norm) |- (records);
\draw[flow] (whisper) -- (records);
\draw[flow] (mel) -- (records);
\draw[flow] (muq) |- (records);
\draw[flow] (records) -- (validation);
\draw[flow] (validation) -- (self);
\draw[flow] (validation) -- (distill);
\draw[optflow] (validation) -- (kaggle);
\draw[optflow] (kaggle) |- (sample);
\draw[flow] (self) -- (sample);
\draw[flow] (distill) -- (sample);
\draw[flow] (sample) -- (output);
\end{tikzpicture}
\caption{Kiến trúc tổng thể của hệ thống sinh nhạc từ văn bản}
\label{fig:pipeline}
\end{figure}

Ở lớp dữ liệu và biểu diễn, đầu vào âm thanh được tách thành vocal và backing. Cả hai được
cộng lại (công thức \eqref{eq:reconstruct_full_mix}) để tạo mel mục tiêu $x_1$ (hỗn hợp cả
bài) mà model học sinh ra. Whisper tạo nội dung lời cùng các đoạn thời gian, trong khi
MuQ-MuLan tạo một vector phong cách toàn cục. Các đầu ra này được quy về cùng đơn vị thời
gian và lưu trong một record thống nhất. Một mẫu huấn luyện vì vậy gồm bộ $(x_1^{\mathrm{voc}},
x_1^{\mathrm{acc}}, l, s)$ — nhưng chỉ $l$ (lời) và $s$ (style) được đưa vào mạng $v_\theta$
làm điều kiện; $x_1^{\mathrm{acc}}$ chỉ tham gia lúc xây dựng mục tiêu, không phải một tensor
điều kiện của model.

Trong lớp học trường vận tốc, sau validation, cùng một dataset hỗ trợ hai nhánh. Nhánh tự huấn luyện tối ưu
MicroDiT trực tiếp bằng CFM ground-truth velocity. Nhánh chưng cất gọi
teacher DiffRhythm2 tại cùng trạng thái nhiễu và timestep, sau đó ánh xạ đầu ra teacher
từ 64 lên 100 mel bins để tạo mục tiêu bổ sung. Hai nhánh dùng chung kiến trúc student
và chỉ khác hàm mất mát, nhờ đó có thể so sánh baseline với distilled student mà không
thay đổi biểu diễn đầu vào.

Ở lớp suy luận và giải mã, hệ thống khởi tạo Gaussian noise trong không gian mel, gọi MicroDiT
nhiều lần để tích phân ODE bằng Euler, rồi đưa mel cuối vào Vocos. Kaggle chỉ là backend
thực thi tùy chọn; nó không thay đổi kiến trúc mô hình hoặc định dạng checkpoint. Cách
phân lớp này giúp cô lập lỗi: có thể kiểm tra riêng record, vòng lặp vocoder, CFM loss,
teacher contract và waveform đầu ra trước khi đánh giá chất lượng âm nhạc tổng thể.

MicroDiT là backbone \textbf{duy nhất} học trường
vận tốc trong hệ thống. Một biến thể thử nghiệm khác (ca từ và mel dùng chung một chuỗi
self-attention thay vì cross-attention riêng, âm vị nhúng bằng một bảng embedding học
từ đầu thay vì XPhoneBERT pretrained) không cải thiện $\mathcal{L}_{\mathrm{gt}}$ nhưng chậm hơn
$\approx4{,}4$ lần mỗi step, nên không được giữ lại trong kiến trúc cuối cùng. Mô hình nhận
noisy mel (hoặc noisy latent, \S\ref{sec:latent_space}) tại thời điểm $t$ cùng
hai nhóm điều kiện: biểu diễn ca từ và neo phong cách toàn cục —
\textbf{không có} tensor backing nào được đưa vào model, dù backing vẫn quyết định mục
tiêu huấn luyện (công thức \eqref{eq:reconstruct_full_mix}). Hình~\ref{fig:microdit_architecture}
tóm tắt luồng xử lý.

Ca từ được chuyển thành âm vị (G2P tiếng Việt có dấu thanh), mã hoá bằng XPhoneBERT
đã đóng băng \cite{xphonebert} cộng một lớp self-attention nhỏ có thể huấn luyện, rồi
được dùng làm điều kiện cho mel thông qua cross-attention ở mỗi khối Transformer — thay
cho một thiết kế "prepend" (nối text và mel vào cùng một chuỗi self-attention) từng được
thử trước đó. Lựa chọn cross-attention theo kết quả của SongGen (arXiv:2502.13128, FAD
1,73 so với 3,56, PER 43,34 so với 56,21 khi so cross-attention với prepend). Timestep và
style embedding (từ MuQ-MuLan) được kết hợp thành một vector điều kiện chung, vừa cộng
trực tiếp vào mel ở đầu vào vừa điều khiển một phép chuẩn hoá thích ứng (adaptive
normalization) trước tầng dự đoán cuối. Từ trạng thái ẩn cuối cùng, hai đầu tuyến tính độc
lập dự đoán velocity của hỗn hợp cả bài (mục tiêu chính) và velocity riêng của vocal (đầu
phụ trợ, chỉ dùng lúc huấn luyện) — mục đích của đầu phụ trợ là tránh model "lãng quên"
tín hiệu vocal, vốn yếu và thưa hơn nhạc đệm, khi tối ưu trực tiếp trên hỗn hợp.

\begin{figure}[H]
\centering
\begin{tikzpicture}[
    >=Latex,
    input/.style={draw, rounded corners=2pt, fill=gray!10,
        minimum height=0.75cm, text width=2.4cm, align=center, font=\scriptsize},
    encoder/.style={draw, rounded corners=2pt, fill=blue!7,
        minimum height=0.85cm, text width=2.9cm, align=center, font=\scriptsize},
    core/.style={draw, rounded corners=2pt, fill=orange!12,
        minimum height=0.9cm, text width=4.3cm, align=center, font=\scriptsize},
    output/.style={draw, rounded corners=2pt, fill=green!9,
        minimum height=0.85cm, text width=3.4cm, align=center, font=\scriptsize},
    flow/.style={-{Latex[length=2mm]}, line width=0.55pt}
]
\node[input] (lyrics) at (-5.1,0) {Ca từ của crop};
\node[input] (mel) at (-1.2,0) {$x_t$\\$B\times T\times100$};
\node[input] (time) at (1.7,0) {Timestep $t$};
\node[input] (style) at (5.1,0) {MuQ-MuLan style\\$B\times512$};

\node[encoder] (g2p) at (-5.1,-1.4)
    {Bộ chuyển tự vị--âm vị (G2P)\\ca từ $\rightarrow$ âm vị IPA};
\node[encoder] (xphonebert) at (-5.1,-2.9)
    {XPhoneBERT đóng băng + MLP\\$E_l\in\mathbb{R}^{B\times L\times D}$};
\node[encoder] (melready) at (-1.2,-2.9)
    {Chiếu tuyến tính $100\rightarrow D$\\$B\times T\times D$};
\node[encoder] (tembed) at (1.7,-2.9)
    {Sinusoidal embedding\\MLP $D\rightarrow D$};
\node[encoder] (sembed) at (5.1,-2.9)
    {Bộ mã hoá phong cách + MLP\\$512\rightarrow D\rightarrow D$};

\node[core] (merge) at (0,-4.45)
    {Khối nhúng đầu vào\\$x_{\mathrm{proj}}+\operatorname{Expand}(E_t)+\operatorname{Expand}(E_s)$, chiếu lại $D\rightarrow D$};
\node[core] (blocks) at (0,-5.95)
    {$K\times$ khối Transformer\\self-attn (mel, RoPE) + cross-attn (mel$\rightarrow$text) + FFN};
\node[core] (adaln) at (0,-7.35)
    {Chuẩn hoá thích ứng cuối\\scale/shift từ $c=E_t+E_s$};
\node[output] (velocity) at (-2.2,-8.7)
    {Chiếu tuyến tính $D\rightarrow100$\\velocity cả bài $\hat{u}_t$};
\node[output] (vocalaux) at (2.2,-8.7)
    {Chiếu tuyến tính phụ trợ $D\rightarrow100$\\velocity vocal (phụ trợ)};

\draw[flow] (lyrics) -- (g2p);
\draw[flow] (g2p) -- (xphonebert);
\draw[flow] (mel) -- (melready);
\draw[flow] (time) -- (tembed);
\draw[flow] (style) -- (sembed);
\draw[flow] (melready) -- (merge);
\draw[flow] (tembed) -- (merge);
\draw[flow] (sembed) -- (merge);
\draw[flow] (merge) -- (blocks);
\draw[flow] (xphonebert) -- (blocks);
\draw[flow] (blocks) -- (adaln);
\draw[flow] (adaln) -- (velocity);
\draw[flow] (adaln) -- (vocalaux);
\end{tikzpicture}
\caption{Kiến trúc và luồng tensor của MicroDiT — điều kiện lời qua cross-attention, không có nhánh backing}
\label{fig:microdit_architecture}
\end{figure}

\subsection{Hàm mất mát}
\label{sec:cfm_training_loss}

Với mỗi batch, thuật toán sinh $x_0\sim\mathcal{N}(0,I)$ cùng kích thước $x_1$ (công thức
\eqref{eq:reconstruct_full_mix}) và lấy một timestep $t\sim\mathcal{U}(0,1)$ độc lập cho
từng mẫu:
\begin{align}
 x_t&=(1-t)x_0+t x_1,\qquad u_t=x_1-x_0,\qquad \hat{u}_t=v_\theta(x_t,t,l,s).
\end{align}
Thay vì MSE trần, hàm mất mát kết hợp bốn thành phần để chống hiện tượng
\emph{regression-to-the-mean}/"distributional averaging" -- xu hướng model hội tụ về một
mức năng lượng trung bình phẳng thay vì phân biệt đúng các đoạn có/không có tiếng hát,
quan sát được khi chỉ dùng MSE trần ở các thử nghiệm ban đầu của đề tài:
\begin{align}
 w &= 1+2\cdot\sigma\!\big(2(\bar{x}_1-Q_{0.55}(\bar{x}_1))\big) \quad\text{(trọng số theo năng lượng frame, chuẩn hoá về trung bình 1)},\\
 \mathcal{L}_{\mathrm{vel}} &= \operatorname{mean}\!\big[w\cdot(\hat{u}_t-u_t)^2\big],\\
 \hat{x}_1 &= x_t+(1-t)\hat{u}_t \quad\text{(mẫu sạch tái tạo một bước)},\\
 \mathcal{L}_{\mathrm{recon}} &= \operatorname{mean}\!\big[w\cdot|\hat{x}_1-x_1|\big],\\
 \mathcal{L}_{\Delta} &= \operatorname{L1}(\Delta_t\hat{x}_1,\Delta_t x_1)+\operatorname{L1}(\Delta_f\hat{x}_1,\Delta_f x_1),\\
 \mathcal{L}_{\mathrm{gt}} &= \mathcal{L}_{\mathrm{vel}}+0{,}15\,\mathcal{L}_{\mathrm{recon}}+0{,}05\,\mathcal{L}_{\Delta}.
 \label{eq:loss_gt}
\end{align}
Nếu trọng số vocal-auxiliary $\lambda_{\mathrm{voc}}>0$, cùng công thức trọng-số-frame được áp lại cho đầu phụ trợ
với mục tiêu $x_1^{\mathrm{voc}}-x_0$, cộng thêm vào loss theo trọng số đó. Khi có ca từ
khác nhau trong cùng batch, một cặp lời "khớp/lệch" được tạo (mỗi mẫu đổi sang lời của một
mẫu khác) và hai số hạng bổ sung được cộng vào: một \emph{contrastive hinge} phạt khi sai
số dự đoán với lời khớp không thấp hơn hẳn sai số với lời lệch, và một số hạng
\emph{sensitivity} phạt khi đổi lời không làm thay đổi đủ nhiều dự đoán (đo bằng RMS
tương đối của hiệu hai velocity); cả hai số hạng đóng vai trò regularizer trong loss và
là cơ sở của tiêu chí chọn checkpoint tốt nhất mô tả ở Chương~\ref{chapter:experiment}
(\S\ref{sec:early_stopping_full}), cùng optimizer, learning-rate schedule và các
hyperparameter huấn luyện cụ thể khác.

\subsection{Suy luận và tổng hợp âm thanh}

Suy luận bắt đầu bằng Gaussian noise $x_0$ có hình dạng $1\times F\times T$ ($F=100$ ở mel-space, $F=64$ ở latent-space, \S\ref{sec:latent_space}). Tại mỗi $t_k$, MicroDiT dự đoán velocity và Euler solver cập nhật trạng thái (mặc định 32 bước, có thể điều chỉnh). Sau bước cuối, mel được giải mã trực tiếp bằng Vocos (hoặc bằng decoder BigVGAN thật nếu đang ở không gian latent). Nếu không tải được Vocos, hệ thống chuyển sang Griffin--Lim 64 vòng lặp và ghi rõ backend đã dùng.

Với bài dài hơn một crop, lời được chia theo segment và các đoạn mel được sinh tuần tự rồi ghép trên trục thời gian. Người dùng có thể chỉ định một bài tham chiếu đã tiền xử lý để lấy style embedding thật, thay vì fallback về vector text-pooled. Checkpoint loại các trọng số XPhoneBERT đã đóng băng vì chúng có thể tải lại từ mô hình gốc, làm giảm mạnh kích thước tệp mà không thay đổi tham số đã học.

\section{Không gian latent thật của DiffRhythm2}
\label{sec:latent_space}

DiffRhythm2 không hoạt động trên mel-spectrogram: teacher thật chạy trên latent 64 chiều
của một Music VAE riêng, ở tần số \textbf{5~Hz} — thấp hơn $\approx19$ lần tần số mel
93,75~Hz mà student vẫn dùng. Phần \S\ref{sec:distill_teacher_call} chỉ sửa việc này ở
\emph{phía truy vấn teacher} lúc chưng cất; hướng này cho \textbf{student chính nó} một
không gian latent nén thật, độc lập với việc có chưng cất hay không.

\subsection{Kiến trúc encoder}
\label{sec:latent_encoder_arch}

DiffRhythm2 công bố \textbf{decoder} (BigVGAN, đã huấn luyện sẵn và công khai) nhưng không
công bố encoder tương ứng. Thay vì train lại một VAE đầy đủ đúng-như-paper (có
discriminator đối kháng, quá tốn/rủi ro cho ngân sách của đề tài), phương pháp giữ decoder
\textbf{đông lạnh} (đã pretrained, chính là decoder thật của teacher) và chỉ train một
encoder tích chập nhỏ mới (khoảng 11 triệu tham số), với tỉ lệ nén thời gian được ràng buộc
khớp chính xác thông số encode-side thật của teacher.

Encoder train trực tiếp trên waveform 24~kHz gốc (vocal+backing đã tách qua Demucs), không
qua bước trung gian mel/Vocos nào — tránh mất pha do một vòng vocoder phụ không cần thiết.
Loss reconstruction là trung bình L1-trên-log-mel qua nhiều thang phân giải STFT, so giữa
audio decoder tái tạo và audio thật — không có số hạng đối kháng (rủi ro bất ổn GAN
training, cố ý để ngoài phạm vi đề tài).

\subsection{Bottleneck xác suất (VAE) và lịch trình KL}
\label{sec:vae_bottleneck}

Bộ mã hoá latent áp dụng reparameterization trick chuẩn của VAE:
$z=\mu+\sigma\epsilon$ lúc train (backprop qua sample thật), $z=\mu$ lúc eval (hành vi xác
định cho suy luận), với hàm mất mát tổng
\begin{equation}
 \mathcal{L} = \mathcal{L}_{\text{recon}} + \beta\cdot\mathcal{L}_{\mathrm{KL}},
 \quad \mathcal{L}_{\mathrm{KL}}=-\tfrac12\,\mathrm{mean}\big(1+\log\sigma^2-\mu^2-\sigma^2\big),
\end{equation}
trong đó $\beta$ là trọng số của số hạng KL. Trọng số này được áp dụng theo lịch trình \textbf{cyclical annealing} (Fu
et al., NAACL 2019) thay vì hằng số: tổng số step huấn luyện chia thành 4 chu kỳ, mỗi chu
kỳ ramp tuyến tính $0\to\beta$ ở nửa đầu rồi giữ nguyên ở nửa sau. Động cơ và
kết quả ablation cho lựa chọn thiết kế này được trình bày ở Thí nghiệm 5
(\S\ref{sec:exp5_vae}), decoder BigVGAN được giữ đông lạnh như mô tả ở
\S\ref{sec:latent_encoder_arch}.

\section{DiffRhythm2 Distillation}
\label{sec:distillation_chapter}

\subsection{Khôi phục đúng hợp đồng teacher}
\label{sec:distill_teacher_call}

Pipeline chưng cất gọi teacher DiffRhythm2 đúng theo cấu hình và giao thức chính thức đi
kèm checkpoint công khai của nó (tokenizer, định dạng token, quy ước timestep), thay vì một
phiên bản đơn giản hoá tự giả định. Teacher luôn ở trạng thái đóng băng, chỉ đóng vai trò
sinh mục tiêu. Student luôn là MicroDiT (backbone duy nhất trong kiến trúc hiện tại,
\S\ref{sec:microdit_architecture}), chạy được trên cả mel-space và latent-space
(\S\ref{sec:latent_space}).

\subsection{Cầu nối tần số 93,75~Hz--5~Hz và 100--64 mel bins}

Khi student chạy ở mel-space, một lệch biểu diễn nghiêm trọng xuất hiện: teacher chạy ở
5~Hz/64 chiều nhưng student ở 93,75~Hz/100 chiều — đưa nguyên chuỗi mel-space vào teacher
khiến attention của nó chạy ngoài phân phối huấn luyện (lệch tần số $\approx19$ lần). Cầu
nối giải quyết lệch này bằng cách resample trục thời gian theo cả hai chiều quanh lượt gọi
teacher, khớp lại đúng mặt nạ attention tự hồi quy theo khối mà teacher được huấn luyện với
(thay vì để nó thấy toàn chuỗi không hạn chế), và một lớp tuyến tính có thể học ánh xạ giữa
64 và 100 chiều để bù lệch số kênh, được huấn luyện đồng thời với phần còn lại của student.

\textbf{Toàn bộ cầu nối này được bỏ qua khi student chạy trên không gian latent.} Nếu
student đã chạy trên latent 64-chiều/5~Hz thật (\S\ref{sec:latent_space}) — đúng tần số và
đúng số chiều của teacher — hệ thống tự động bỏ qua cả bước ánh xạ kênh lẫn bước resample
thời gian, vì hai phía đã cùng một biểu diễn ngay từ đầu; áp lại công thức resample dành
cho mel-space lên một chuỗi đã ở 5~Hz sẽ làm hỏng nó thêm một lần nữa, không phải chỉnh
đúng. Khoảng cách phân phối 18,75 lần giữa hai tần số vì vậy không còn tồn tại.

\subsection{Hàm mất mát hỗn hợp}
\label{sec:distill_mixed_loss}

Student được tối ưu bởi:
\begin{equation}
 \mathcal{L}=(1-\alpha)
 \operatorname{L1}(\hat{u}_t,\tilde{v}_T)
 +\alpha\,\mathcal{L}_{\mathrm{gt}}
\end{equation}
với $\mathcal{L}_{\mathrm{gt}}$ theo đúng công thức \eqref{eq:loss_gt} (không phải MSE
trần). Số hạng teacher-matching dùng \textbf{L1} chứ không phải MSE, theo khuyến nghị
literature (Dieleman 2024; DMD/ADM) để tránh xu hướng "làm mượt/lấy trung bình" của MSE
trong distillation cho model sinh. Cấu hình mặc định trong mã nguồn là $\alpha=0{,}5$,
nhưng đo tỉ lệ khung có giọng trên 6 bài khác nhau cho ba giá trị $\alpha$ -- $60{,}7\%$
ở $\alpha=0{,}5$, $\mathbf{92{,}4\%}$ ở $\alpha=0{,}8$, và $77{,}8\%$ ở $\alpha=0{,}9$ (vocal
thật: $74{,}3\%$) -- xác nhận $\alpha\approx0{,}8$ là điểm tối ưu thật (quan hệ phi tuyến,
không phải "$\alpha$ cao hơn luôn tốt hơn", và không phải nhiễu ngẫu nhiên từ một checkpoint
đơn lẻ) — nên dùng $\alpha=0{,}8$ làm mặc định thực tế, không phải $0{,}5$. Vocal-aux loss và một loss REPA
tuỳ chọn (mặc định tắt) được cộng thêm vào không điều kiện nếu trọng số tương ứng $>0$.
Nếu teacher hoặc tokenizer không tải được, $\alpha$ bị ép về $1{,}0$ (chỉ còn
$\mathcal{L}_{\mathrm{gt}}$) và trạng thái này được báo cáo rõ ràng trong log, không âm thầm.

\subsection*{Tóm tắt lựa chọn thiết kế}

Phương pháp được xây dựng quanh bốn nguyên tắc. Thứ nhất, biểu diễn trung gian phải nhất
quán từ tiền xử lý tới vocoder, vì sai hợp đồng mel có thể che lấp hoàn toàn chất lượng mô
hình sinh. Thứ hai, các điều kiện được phân vai: lyrics cho nội dung (qua cross-attention
hoặc chuỗi chung, tuỳ backbone) và MuQ-MuLan cho phong cách toàn cục; backing chỉ tham gia
lúc xây mục tiêu huấn luyện (hỗn hợp cả bài), không phải một tensor điều kiện của model.
Thứ ba, chưng cất phải được kiểm chứng ở mức đường truyền gradient và hợp đồng teacher
trước khi bàn tới lợi ích chất lượng. Thứ tư, khi cho student một không gian biểu diễn mới
(latent thật của teacher), phải tự kiểm chứng không gian đó bằng cách giải mã trực tiếp
latent ground-truth trước khi tin bất kỳ kết quả CFM training nào trên nó — một encoder
suy biến có thể trông ổn về mặt kỹ thuật (không NaN, loss giảm) mà vẫn không mang thông tin
âm nhạc nào. Các nguyên tắc này quyết định thiết kế thí nghiệm ở chương tiếp theo.


\chapter{THỰC NGHIỆM}
\label{chapter:experiment}

\section{Mục tiêu và câu hỏi nghiên cứu}

Thực nghiệm được thiết kế theo bốn tầng. Tầng thứ nhất xác nhận tính đúng của biểu diễn âm thanh và vocoder. Tầng thứ hai xác nhận pipeline dữ liệu--huấn luyện--suy luận có thể chạy đầu cuối trên dữ liệu tiếng Việt. Tầng thứ ba xác nhận teacher DiffRhythm2 thật đã tham gia vào gradient chưng cất và đo chi phí của cấu hình này. Tầng thứ tư so sánh không gian biểu diễn của student (mel 100-chiều/93,75~Hz so với latent thật 64-chiều/5~Hz của teacher). Cách tổ chức theo tầng giúp kết luận tương ứng với bằng chứng, thay vì dùng một file âm thanh đầu ra để suy diễn cho mọi thành phần.

Bốn câu hỏi nghiên cứu được đặt ra:
\begin{description}
 \item[RQ1:] Hợp đồng mel 100 bins với Vocos 24~kHz có tái tạo waveform ổn định hơn các đường giải mã không tương thích hay không?
 \item[RQ2:] Pipeline có xử lý được dữ liệu bài hát tiếng Việt, huấn luyện MicroDiT bằng CFM và sinh waveform hữu hạn, không rỗng hay không?
 \item[RQ3:] Quy trình distillation có thực sự gọi teacher, ánh xạ 64--100 mel bins và cập nhật student nhỏ hay không; kết quả hiện có cho phép kết luận gì về lợi ích của teacher?
 \item[RQ4:] Cho student không gian latent thật của teacher (thay vì mel) có thay đổi được vấn đề "âm thanh sinh ra nghe như nhiễu" đã lặp lại ở các thí nghiệm mel-space trước đó không, và nguyên nhân của sự khác biệt (nếu có) là gì?
 \item[RQ5:] Bottleneck xác suất (VAE, mu/logvar + KL divergence) có khắc phục được hiện tượng collapse của latent encoder ở quy mô dữ liệu lớn hay không, và trọng số $\beta$ nào là tối ưu?
\end{description}

\section{Nguồn dữ liệu và nguyên tắc báo cáo}

Báo cáo chỉ sử dụng các giá trị có trong log, artifact hoặc checkpoint của dự án. Dữ liệu đánh giá gồm ba nhóm. Nhóm kiểm tra vocoder dùng các đoạn âm thanh thật để thực hiện vòng lặp waveform $\rightarrow$ mel $\rightarrow$ waveform. Nhóm đầu cuối trên Kaggle gồm 12 bài và hoàn thành tiền xử lý 12/12. Nhóm cục bộ gồm 2 bài cho smoke test đầu cuối và một thí nghiệm distillation 30 epoch. Một tập ứng viên lớn hơn có 201 bài, dung lượng khoảng 2,52~GB, được dùng để ước lượng chi phí tiền xử lý chứ chưa được dùng để công bố kết quả huấn luyện hoàn chỉnh.

Cần phân biệt reconstruction và generation. Trong reconstruction, mel được trích từ chính waveform tham chiếu nên correlation và RMSE có cặp mẫu tương ứng. Trong generation, waveform mới không có waveform đích duy nhất; do đó correlation với bài gốc không phải thước đo chất lượng sinh hợp lệ. Thực nghiệm generation chỉ báo cáo sanity metrics và loss, còn đánh giá âm nhạc đầy đủ cần listening test hoặc metric ngữ nghĩa riêng.

\section{Cấu hình thí nghiệm}

\subsection{Cấu hình baseline và distillation}

Baseline dùng MicroDiT $D=256$, 4 transformer layers, 4 heads; text encoder (XPhoneBERT) bị đóng băng. CFM dùng crop 128 frames và AdamW trên các tham số trainable. Lần chạy Kaggle huấn luyện 40 epoch, mỗi epoch 3 step, tổng cộng 120 optimization steps. Đây là phép kiểm tra khả năng học và tích hợp, chưa phải lịch huấn luyện hội tụ cho mô hình sinh nhạc quy mô lớn.

Thí nghiệm distillation cục bộ dùng student nhỏ $D=128$, 2 layers, batch size 2 và 30 epoch trên 2 bài. Teacher DiffRhythm2 được tải từ checkpoint thật, có 1.136.249.664 tham số; student có 745.188 tham số trainable trong cấu hình ghi nhận. Hệ số trộn là $\alpha=0{,}5$. Ngoài cấu hình này, project định nghĩa các biến thể $\alpha\in\{0{,}2;0{,}5;0{,}8\}$ và baseline/student nhỏ để phục vụ đối chứng khi có đủ tài nguyên.

\begin{table}[H]
\centering
\caption{Các cấu hình thực nghiệm đã ghi nhận}
\label{tab:experiment_config}
\small
\begin{tabularx}{\textwidth}{p{2.8cm}p{2.8cm}p{2.5cm}X}
\toprule
\textbf{Thí nghiệm} & \textbf{Dữ liệu} & \textbf{Mô hình} & \textbf{Mục tiêu} \\
\midrule
Vocoder loopback & Đoạn âm thanh thật & Vocos / Griffin--Lim & Kiểm tra hợp đồng feature--decoder. \\
Kaggle baseline & 12 bài & $D=256$, 4 layers & Xác nhận pipeline đầu cuối và CFM training. \\
Local baseline & 2 bài & MicroDiT mặc định & Smoke test trên môi trường cục bộ. \\
Local distillation & 2 bài & $D=128$, 2 layers & Xác nhận teacher thật, adapter và gradient. \\
\bottomrule
\end{tabularx}
\end{table}

\subsection{Cấu hình huấn luyện, chọn checkpoint và dừng sớm}
\label{sec:early_stopping_full}

Toàn bộ tham số không đóng băng của CFM/distillation (\S\ref{sec:cfm_training_loss})
được tối ưu bằng AdamW, cùng trung bình trượt theo hàm mũ (EMA, hệ số suy giảm $0{,}999$),
lịch trình learning rate khởi động rồi suy giảm theo hàm cosine, huấn luyện độ chính xác
hỗn hợp trên GPU và cắt ngưỡng gradient (chuẩn $1,0$). Trước khi train, các bản ghi hữu
dụng được chia train/validation bằng một hàm băm mật mã quyết định theo seed và mã số bản
ghi — không dùng số ngẫu nhiên, nên cách chia ổn định qua các lần resume/train lại; tỉ lệ
validation mặc định $5\%$, tối đa 128 bản ghi.

Sau mỗi epoch, một checkpoint chỉ được coi là "tốt nhất" nếu \textbf{đồng thời}: (a)
validation loss (đánh giá bằng trọng số EMA, seed cố định) giảm hơn một ngưỡng tối thiểu;
\textbf{và} (b) chỉ số sensitivity (\S\ref{sec:cfm_training_loss}) không thấp hơn một sàn
tối thiểu — một checkpoint cải thiện validation loss bằng cách lờ hẳn ca từ sẽ không qua
được điều kiện (b). Việc train dừng sớm khi đã đạt tối thiểu 8 epoch \textbf{và} không cải
thiện được checkpoint tốt nhất qua 4 epoch liên tiếp — nghĩa là số epoch truyền vào chỉ là
một mức trần, cơ chế này mới là thứ quyết định khi nào dừng thật.

\subsection{Chỉ số đánh giá}

Với waveform tham chiếu $y$ và waveform tái tạo $\hat{y}$ đã căn cùng độ dài, hệ số tương quan Pearson đo độ giống về hình dạng:
\begin{equation}
 r=\frac{\sum_n(y_n-\bar{y})(\hat{y}_n-\bar{\hat{y}})}
 {\sqrt{\sum_n(y_n-\bar{y})^2}\sqrt{\sum_n(\hat{y}_n-\bar{\hat{y}})^2}}.
\end{equation}
RMSE đo sai số biên độ:
\begin{equation}
 \operatorname{RMSE}=\sqrt{\frac{1}{N}\sum_{n=1}^{N}(y_n-\hat{y}_n)^2}.
\end{equation}
Các giá trị RMSE trong bảng vocoder được giữ theo đúng thang đo của script đánh giá tại thời điểm chạy; do các phiên bản pipeline có chuẩn hóa khác nhau, chỉ nên so sánh trong cùng nhóm loopback.

Với waveform sinh, các sanity metrics gồm peak amplitude, RMS, tỷ lệ mẫu gần không và kiểm tra NaN/Inf. Chúng trả lời tín hiệu có hữu hạn, có năng lượng và không bị im lặng hay clipping nghiêm trọng hay không; chúng không đo độ hay, tính đúng lời hoặc chất lượng giai điệu. Với mô hình, báo cáo CFM ground-truth loss, teacher-matching loss, số tham số và thời gian forward teacher.

\section{Thí nghiệm 1: kiểm chứng vocoder}
\label{sec:e1_vocoder}

Các lần thử ban đầu dùng mel tự xây dựng hoặc thay đổi kích thước bằng nội suy trước khi đưa vào Vocos. Mặc dù tensor đúng 100 bins, kết quả không duy trì tốt waveform vì miền đặc trưng không khớp checkpoint. Sau khi dùng feature extractor gốc của Vocos trong cả tiền xử lý và loopback, correlation tăng rõ rệt.

\begin{table}[H]
\centering
\caption{Kết quả vòng lặp waveform--mel--waveform}
\label{tab:vocoder_results}
\small
\begin{tabularx}{\textwidth}{Xrr}
\toprule
\textbf{Đường xử lý} & \textbf{Correlation} & \textbf{RMSE} \\
\midrule
Mel cũ, Griffin--Lim với pha khởi tạo không phù hợp & 0,149 & 4,37 \\
Mel cũ, nội suy bilinear rồi Vocos & 0,686 & 2,92 \\
Mel cũ, Griffin--Lim đúng cấu hình phổ & 0,495 & 8,12 \\
Feature extractor Vocos, phiên bản tích hợp ban đầu & 0,896 & 3,11 \\
Feature extractor và decode Vocos đã đồng bộ & \textbf{0,997} & \textbf{1,78} \\
Mel đồng bộ, Griffin--Lim 64 vòng & 0,960 & 2,85 \\
\bottomrule
\end{tabularx}
\end{table}

Kết quả tốt nhất đạt correlation 0,997, trong khi Griffin--Lim trên mel đã đồng bộ đạt 0,960. Một lần kiểm tra trên Kaggle đạt 0,993; một bài khác chạy cục bộ đạt 0,986. Các giá trị lặp lại ở nhiều môi trường cho thấy đường Vocos không chỉ hoạt động trên một file duy nhất.

So sánh này trả lời RQ1 theo hai ý. Thứ nhất, số mel bins chỉ là một phần của hợp đồng; nội suy một mel không tương thích lên 100 bins không khôi phục được thống kê mà Vocos đã học. Thứ hai, khi feature extractor và decoder đồng bộ, neural vocoder bảo toàn waveform tốt hơn đường cũ. Vì vậy, các thí nghiệm generation về sau sử dụng biểu diễn Vocos-native; nếu chuyển sang Griffin--Lim, backend được ghi riêng để tránh trộn kết quả.

\section{Thí nghiệm 2: pipeline đầu cuối}

\subsection{Kết quả tiền xử lý và huấn luyện}

Trong lần chạy Kaggle hoàn chỉnh, 12/12 bài đi qua pipeline tiền xử lý và tạo record hợp lệ. Baseline được huấn luyện 40 epoch, 3 step mỗi epoch. CFM loss cuối cùng là 5,34; toàn bộ 120 step hoàn thành trong khoảng 11 giây ở môi trường chạy đã ghi nhận. Loss hữu hạn xác nhận các tensor mel, text, style, timestep và backing có thể đi qua mô hình và backward mà không phát sinh lỗi số.

\begin{table}[H]
\centering
\caption{Tóm tắt kiểm tra đầu cuối}
\label{tab:e2e_results}
\small
\begin{tabularx}{\textwidth}{p{3.2cm}p{2.7cm}X}
\toprule
\textbf{Thiết lập} & \textbf{Kết quả} & \textbf{Sanity check waveform} \\
\midrule
Kaggle, 12 bài & 12/12 tiền xử lý; 120 step; loss cuối 5,34 & Peak khoảng 0,8; RMS 0,08--0,15; silence $<0,13\%$; không NaN/Inf. \\
Cục bộ, 2 bài & 2/2 tiền xử lý; 5 epoch, 1 step/epoch & RMS 0,07--0,10; silence $<0,2\%$; không NaN/Inf; loopback Vocos 0,986. \\
\bottomrule
\end{tabularx}
\end{table}

Waveform sinh có năng lượng khác không, tỷ lệ silence thấp và không chứa NaN/Inf. Peak khoảng 0,8 phản ánh bước chuẩn hóa đầu ra, không phải bằng chứng tự thân về chất lượng. Kết quả này trả lời phần tính khả thi của RQ2: pipeline thực sự tạo được dữ liệu, cập nhật mô hình và sinh waveform nghe được về mặt tín hiệu.

\subsection{Phân tích ý nghĩa kết quả}

CFM loss 5,34 không thể so sánh trực tiếp với các công trình khác khi khác preprocessing, normalization và kích thước tensor. Giá trị của nó ở đây là chỉ báo nội bộ rằng phép hồi quy velocity ổn định qua 120 step. Tương tự, RMS và silence ratio loại trừ hai lỗi nghiêm trọng là đầu ra rỗng và bùng nổ số, nhưng chưa trả lời lời hát có đúng hay giai điệu có tự nhiên.

Lịch huấn luyện ngắn cũng khiến model có xu hướng học các thống kê cục bộ của tập 12 bài. Vì vậy, báo cáo không dùng thí nghiệm này để tuyên bố khả năng tổng quát hóa. Đóng góp của nó là kiểm chứng toàn bộ chuỗi tích hợp trên dữ liệu thật, tạo nền tảng để mở rộng lịch train và đánh giá người nghe trong giai đoạn sau.

\section{Thí nghiệm 3: chưng cất DiffRhythm2}

\subsection{Kiểm chứng cơ chế}

Log thí nghiệm xác nhận chế độ chưng cất thực sự được kích hoạt; checkpoint teacher và tokenizer thật đã được tải. Teacher có 1.136.249.664 tham số, trong khi student nhỏ có 745.188 tham số trainable, tương ứng tỷ lệ xấp xỉ 1.525 lần theo số tham số. Trên CPU, một forward teacher với batch size 1 mất khoảng 3 giây. Sự chênh lệch cho thấy động cơ nén là rõ ràng, đồng thời giải thích vì sao distillation trở thành nút thắt thời gian.

Lớp ánh xạ 64--100 mel bins nằm trong optimizer cùng student. Teacher forward được đóng băng, nhưng đầu ra teacher đi qua lớp ánh xạ này bên ngoài phạm vi đóng băng đó nên vẫn nhận gradient. Việc log riêng teacher loss và ground-truth loss xác nhận tổng loss thực sự gồm hai nguồn. Như vậy, RQ3 được trả lời ở mức cơ chế: đây là một lần chưng cất thật, không phải baseline được đổi tên.

\subsection{Kết quả và diễn giải}

Sau 30 epoch trên 2 bài, ground-truth loss cuối được ghi nhận khoảng 17,9 đối với student distillation và 15,7 đối với baseline tương ứng. Loss từng step dao động từ khoảng 3,5 tới 229. Biên độ dao động lớn là có thể dự đoán vì mỗi epoch chỉ có một step, timestep $t$ và crop đều được lấy ngẫu nhiên; một điểm cuối đơn lẻ không đại diện cho đường cong kỳ vọng.

\begin{table}[H]
\centering
\caption{Kết quả thí nghiệm distillation cục bộ}
\label{tab:distill_results}
\begin{tabularx}{\textwidth}{p{4.0cm}p{3.2cm}X}
\toprule
\textbf{Đại lượng} & \textbf{Giá trị ghi nhận} & \textbf{Diễn giải} \\
\midrule
Teacher parameters & 1.136.249.664 & Teacher DiffRhythm2 thật đã tải. \\
Student trainable parameters & 745.188 & Student nhỏ hơn khoảng 1.525 lần. \\
Teacher forward, CPU, batch 1 & Khoảng 3 giây & Chi phí teacher chi phối vòng train cục bộ. \\
Loss GT cuối, distilled / baseline & 17,9 / 15,7 & Không cho thấy cải thiện trong lần chạy nhỏ. \\
Khoảng dao động loss theo step & 3,5--229 & Phương sai cao do 1 step/epoch và lấy mẫu ngẫu nhiên. \\
\bottomrule
\end{tabularx}
\end{table}

Kết quả không chứng minh distillation cải thiện chất lượng, nhưng cũng chưa phải một đối chứng đủ lực để bác bỏ phương pháp. Cỡ mẫu 2 bài, chỉ 30 update và thiếu nhiều seed khiến sai khác 2,2 đơn vị nhỏ hơn độ biến thiên theo step. Ngoài ra, adapter trainable và reference conditioning đã được hoàn thiện sau lần chạy ghi nhận; cần chạy lại ma trận baseline--distilled trên cùng split, cùng seed và nhiều update hơn trước khi kết luận.

\section{Thí nghiệm 4: không gian mel so với không gian latent thật của teacher}
\label{sec:exp4_mel_vs_latent}

\subsection{Thiết lập}

Ba điều kiện được so sánh, đo bằng độ phẳng phổ, tỉ lệ khung có giọng và
SD cao độ, với neo là nhiễu trắng và vocal thật:
\begin{enumerate}[label=(\Alph*)]
 \item \textbf{Mel-space} -- MicroDiT, mốc \emph{exp06} ($\alpha=0{,}8$,
   \S\ref{sec:distill_mixed_loss}), 25 epoch đã hội tụ.
 \item \textbf{Latent-space, encoder chưa ổn định hoá} -- giải mã trực tiếp latent
   \emph{ground-truth} (bỏ qua CFM student) để cô lập chất lượng encoder.
 \item \textbf{Latent-space, encoder đã ổn định hoá + CFM student} -- checkpoint tốt nhất
   hiện có, mới train 13/300 epoch (dừng vì hết quota Kaggle), dùng một backbone thử
   nghiệm không giữ lại trong kiến trúc cuối cùng (\S\ref{sec:microdit_architecture}), không
   phải MicroDiT.
\end{enumerate}

\begin{table}[H]
\centering
\caption{So sánh mel-space và latent-space (số liệu sơ bộ — xem lưu ý bên dưới)}
\label{tab:mel_vs_latent}
\small
\begin{tabularx}{\textwidth}{p{2.6cm}>{\centering\arraybackslash}p{1.5cm}>{\centering\arraybackslash}p{1.5cm}>{\centering\arraybackslash}p{1.6cm}X}
\toprule
\textbf{Điều kiện} & \textbf{Độ phẳng phổ} & \textbf{Tỉ lệ voiced} & \textbf{SD cao độ (bán cung)} & \textbf{Ghi chú} \\
\midrule
Nhiễu trắng (mốc) & 0,562 & 0,040 & 0,49 & tham chiếu \\
Vocal thật (mốc) & 0,056 & 0,743 & 6,39 & tham chiếu \\
(A) Mel-space, exp06 & 0,011 & \textbf{0,924} & 0,91 & tỉ lệ voiced cao nhưng gần như đứng yên một nốt \\
(B) Latent-space, encoder chưa fix & 0,0001--0,0005 & 0,95--0,99 & \textbf{0,83--0,97} & không phải nhiễu trắng, nhưng gần như đứng yên một nốt \\
(C) Latent-space, encoder đã fix + CFM (13/300 epoch) & 0,00095 & 0,21 & \textbf{5,70} & giai điệu thật, gần vocal thật hơn hẳn (A)/(B) \\
\bottomrule
\end{tabularx}
\end{table}

\subsection{Kết quả}

(A) và (B) đều đạt SD cao độ$<1$ semitone -- gần như đứng yên một nốt -- dù
nguyên nhân khác nhau (mel-space: quy mô dữ liệu/step; latent-space: encoder chưa ổn định
hoá, \S\ref{sec:latent_space}). Sau khi ổn định hoá encoder, (C) đạt
SD cao độ$=5{,}70$ ($\approx89\%$ vocal thật, so với $\approx14\%$ của (A)) dù
mới train 13/300 epoch so với 25 epoch đã hội tụ của (A) -- một so sánh không cân bằng
ngân sách nhưng có lợi cho latent-space. Tỉ lệ khung có giọng của (C) ($0{,}21$) thấp
hơn cả (A) và vocal thật, phản ánh checkpoint chưa hội tụ đủ để giữ pitch ổn định xuyên
suốt đoạn có lời. Cho student không gian latent thật của teacher vì vậy cải thiện rõ rệt
SD cao độ, xác nhận thêm bằng nghe thử trực tiếp (2026-07-22) -- lần đầu tiên trong đề tài
một mẫu CFM cho giai điệu nghe được thay vì nhiễu/đơn điệu. Tuy nhiên (C) chưa hội tụ và
dùng backbone đã bị loại khỏi kiến trúc cuối cùng, nên số liệu ở bảng
\ref{tab:mel_vs_latent} là tạm thời; kết quả đầy đủ trên MicroDiT sẽ có sau khi huấn luyện
chưng cất hoàn tất (\S\ref{sec:future_work}).

\section{Thí nghiệm 5: ablation bottleneck xác suất (VAE) của latent encoder}
\label{sec:exp5_vae}

\subsection{Thiết lập}

Encoder ở Thí nghiệm 4 dòng (B) collapse hoàn toàn khi chỉ dùng loss reconstruction thuần
($SD cao độ\approx0{,}44$--$0{,}78$). Bốn cấu hình được so sánh trên
toàn bộ 1839 bài (10 epoch cho hai cấu hình đầu; 5 epoch tiếp tục từ checkpoint có sẵn cho
hai cấu hình sau), đo trên audio thật (không qua CFM student): $\sigma_{\text{mean}}$
(bottleneck có thật sự xác suất), SD cao độ (chất lượng
reconstruction), và khoảng cách $\mu$ giữa các bài $\lVert\mu_A-\mu_B\rVert$ (encoder có
phân biệt được các bài hay bị mean collapse).

\subsection{Kết quả}

\begin{table}[H]
\centering
\caption{Ablation $\beta$ cho bottleneck xác suất của bộ mã hoá latent}
\label{tab:kl_ablation}
\small
\begin{tabularx}{\textwidth}{p{2.6cm}>{\centering\arraybackslash}p{1.6cm}>{\centering\arraybackslash}p{1.4cm}>{\centering\arraybackslash}p{1.5cm}X}
\toprule
\textbf{Cấu hình} & $\boldsymbol{\sigma_{\text{mean}}}$ & \textbf{SD cao độ} & $\boldsymbol{\mu}$\textbf{-distance} & \textbf{Ghi chú} \\
\midrule
Không VAE (2 lần thử) & — & 0,44 / 0,78 & — & collapse hoàn toàn \\
$\beta=10^{-4}$ hằng số & 0,003--0,005 & 5,85$^{*}$ & — & $^{*}$KL vanishing: $\sigma\approx0$, bottleneck vô tác dụng dù SD cao độ trông tốt \\
$\beta=0{,}05$ cyclical & 0,11--0,14 & 4,72 & 0,80--0,91 & VAE thật hoạt động lần đầu \\
$\beta=0{,}15$ cyclical & 0,245--0,273 & \textbf{6,01} & 0,84--0,91 & \textbf{tốt nhất} \\
$\beta=0{,}3$ cyclical & 0,34--0,37 & 3,72 & 0,81--0,88 & vượt ngưỡng, over-regularization \\
\bottomrule
\end{tabularx}
\end{table}
\emph{Mốc tham chiếu: audio thật cùng 5 bài cho SD cao độ trung bình $9{,}46$.
$\mu$-distance đo trên checkpoint có $\sigma$ khoẻ mạnh ($\ge0{,}05$); không đo được ở hai
dòng đầu vì không có $\sigma$ hoặc $\sigma\approx0$ khiến phép đo mất ý nghĩa.}

\subsection{Diễn giải}

Ở $\beta=10^{-4}$ hằng số, SD cao độ$=5{,}85$ trông tốt nhưng
$\sigma_{\text{mean}}\approx0{,}003$--$0{,}005$ cho thấy bottleneck xác suất đã collapse --
hiện tượng \textbf{"KL vanishing"} (Fu et al., NAACL 2019): mạng loại bỏ nhiễu
reparameterization để tối ưu reconstruction thuần tuý, khiến chỉ số trông tốt \emph{vì}
encoder gần như tất định. Lịch trình cyclical annealing khắc phục được điều này, nhưng tồn
tại một điểm tối ưu không đơn điệu: $\sigma$ và SD cao độ cùng tăng từ
$0{,}05\to0{,}15$, nhưng SD cao độ giảm khi tăng tiếp lên $0{,}3$
(over-regularization, Rivera 2023). $\mu$-distance ($0{,}80$--$0{,}91$, xấp xỉ độ lệch
chuẩn nội tại $\approx1{,}0$) không ủng hộ giả thuyết "mean collapse" từng được đề xuất
cho các giá trị $\beta$ nhỏ hơn. Bottleneck xác suất vì vậy là cần thiết để tránh collapse
ở quy mô dữ liệu lớn, nhưng chỉ khi kết hợp lịch trình cyclical annealing thay vì $\beta$
hằng số; checkpoint $\beta=0{,}15$ là cấu hình tốt nhất tìm được và được dùng làm đầu vào
cho bước chưng cất cuối cùng (Chương~\ref{chapter:conclusion}).

\chapter{KẾT LUẬN VÀ HƯỚNG PHÁT TRIỂN}
\label{chapter:conclusion}

\section{Kết luận}

Đề tài đã xây dựng một pipeline sinh vocal/audio tiếng Việt hoàn chỉnh dựa trên
MicroDiT và Conditional Flow Matching, bao phủ toàn bộ vòng đời kỹ thuật: tách
nguồn, nhận dạng lời, trích đặc trưng phong cách, biểu diễn phổ/latent,
huấn luyện, chưng cất tri thức và suy luận. Ba kết luận chính rút ra từ các thí
nghiệm ở Chương~\ref{chapter:experiment}:

Thứ nhất, hợp đồng biểu diễn âm thanh quyết định trực tiếp độ trung thực khi
giải mã (RQ1, \S\ref{sec:e1_vocoder}): một mel-spectrogram đúng cấu hình với
vocoder cho correlation loopback $0{,}997$, so với dưới $0{,}9$ khi có sai lệch
cấu hình -- một điều kiện cần cho mọi thí nghiệm sinh sau đó, dù không phải yếu
tố quyết định chất lượng âm nhạc. Thứ hai, không gian latent thật của teacher (64
chiều, 5~Hz) cho tín hiệu sinh nhạc tốt hơn rõ rệt so với mel-space
(RQ4, \S\ref{sec:exp4_mel_vs_latent}), với điều kiện bộ mã hóa latent tự huấn
luyện phải có bottleneck xác suất (VAE) cùng lịch trình cyclical KL annealing để
tránh collapse ở quy mô dữ liệu lớn (RQ5, \S\ref{sec:exp5_vae}) -- nếu không, cả
hai không gian biểu diễn cho cùng một triệu chứng suy biến (gần như đứng yên một
cao độ), dù nguyên nhân khác nhau. Thứ ba, cơ chế chưng cất tri thức từ teacher
DiffRhythm2 đã được xác nhận hoạt động đúng (RQ3, \S\ref{sec:distillation_chapter}):
teacher thật, lớp ánh xạ 64--100 mel bins và gradient qua hai nguồn loss đều được
kiểm chứng; tuy nhiên bằng chứng hiện có mới đủ xác nhận cơ chế, chưa đủ để khẳng
định mức cải thiện chất lượng cụ thể.

Về mặt triển khai, checkpoint cuối cùng chỉ lưu các trọng số có thể huấn luyện cùng
metadata cấu hình cần thiết để tải lại đúng phụ thuộc, bỏ qua các trọng số nền tảng đã
đóng băng (có thể tải lại từ mô hình gốc) — giảm dung lượng từ khoảng 1,1~GB xuống còn
khoảng 67~MB mà không đổi phép tính của MicroDiT, thuận tiện hơn cho việc lưu phiên bản và
khôi phục thí nghiệm.

Tại thời điểm hoàn thành báo cáo, checkpoint encoder tốt nhất
($\beta=0{,}15$, Thí nghiệm 5) đang được dùng để huấn luyện chưng cất
quy mô đầy đủ trên toàn bộ dữ liệu, nhưng quá trình này chưa hội tụ nên chưa có số
liệu chất lượng cuối cùng. Do đó, kết luận phù hợp nhất là đề tài đã hoàn thành một
pipeline nghiên cứu khả dụng, minh bạch về giới hạn và có cơ sở thực nghiệm để mở
rộng, chứ chưa phải một hệ thống sinh bài hát hoàn chỉnh. Các nút thắt còn lại
(chất lượng chưa ổn định, độ rõ lời chưa cao) không có bằng chứng nào cho thấy là
giới hạn cấu trúc của phương pháp; nhiều khả năng đây là vấn đề về quy mô -- dữ
liệu huấn luyện, số bước/epoch, hoặc tài nguyên tính toán -- và cần được kiểm
chứng bằng cách mở rộng thêm ở các hướng nêu dưới đây.

\section{Hướng phát triển}
\label{sec:future_work}

Dựa trên các kết luận trên, đề tài đề xuất các hướng phát triển tiếp theo theo thứ
tự ưu tiên:
\begin{enumerate}[label=\arabic*)]
    \item Huấn luyện chưng cất trên checkpoint encoder
      $\beta=0{,}15$ tới hội tụ thật, thay vì bị giới hạn bởi thời
      gian/quota GPU như hiện tại; đánh giá bằng cả SD cao độ lẫn chỉ số
      độ rõ lời (nhận dạng tiếng nói tự động) trên nhiều mẫu hơn một vài bài.
    \item Nếu còn dư địa tài nguyên, thử tách riêng số hạng $\mathrm{KL}_G$/$\mathrm{KL}_I$ theo Rivera
      (arXiv:2309.13160) thay cho KL chuẩn, nhằm giảm rủi ro over-regularization đã
      quan sát được ở $\beta=0{,}3$ (\S\ref{sec:exp5_vae}).
    \item Nâng cao số lượng và chất lượng của tập dữ liệu nhạc Việt Nam.
    \item Tổ chức các phương pháp đánh giá chủ quan như MOS/CMOS với người nghe và đánh giá độ rõ lời tiếng Việt.
    \item Mở rộng mô hình, tích hợp các yếu tố khác như thanh điệu, alignment và pitch/F0 hoặc melody.
    \item Kiểm soát các nguồn sai lệch của pipeline dữ liệu trước khi mở rộng quy mô: rò rỉ
      giọng hát vào phần backing do tách nguồn chưa hoàn hảo, và mốc thời gian theo đoạn
      (không chính xác tới từng âm vị) giới hạn khả năng điều khiển lời theo frame.
    \item Chuẩn hoá protocol thực nghiệm khi mở rộng: chia train/validation theo bài hát
      hoặc ca sĩ để tránh rò rỉ dữ liệu, cố định seed và số bước cập nhật, đánh giá trên
      nhiều checkpoint thay vì chỉ dựa vào loss cuối, và kiểm tra quyền sử dụng bản ghi/ca
      từ trước khi phổ biến dữ liệu hoặc audio sinh ra.
\end{enumerate}

% =============================================================
% TÀI LIỆU THAM KHẢO -- đặt trực tiếp để file .tex tự chứa
% =============================================================
\clearpage
\phantomsection
\addcontentsline{toc}{chapter}{TÀI LIỆU THAM KHẢO}
\begin{thebibliography}{99}

\bibitem{musicgen}
J. Copet et al.,
\textit{Simple and Controllable Music Generation},
Advances in Neural Information Processing Systems, 2023.

\bibitem{audioldm}
H. Liu et al.,
\textit{AudioLDM: Text-to-Audio Generation with Latent Diffusion Models},
Proceedings of the International Conference on Machine Learning, 2023.

\bibitem{diffrhythm}
ASLP-lab et al.,
\textit{DiffRhythm: Blazingly Fast and Embarrassingly Simple End-to-End
Full-Length Song Generation with Latent Diffusion},
arXiv:2503.01183, 2025.


\bibitem{mel}
S. S. Stevens, J. Volkmann, and E. B. Newman,
``A Scale for the Measurement of the Psychological Magnitude Pitch,''
\textit{The Journal of the Acoustical Society of America}, vol. 8, no. 3,
pp. 185--190, 1937.

\bibitem{mfcc}
S. B. Davis and P. Mermelstein,
``Comparison of Parametric Representations for Monosyllabic Word Recognition in
Continuously Spoken Sentences,''
\textit{IEEE Transactions on Acoustics, Speech, and Signal Processing}, vol. 28,
no. 4, pp. 357--366, 1980.

\bibitem{griffinlim}
D. Griffin and J. Lim,
``Signal Estimation from Modified Short-Time Fourier Transform,''
\textit{IEEE Transactions on Acoustics, Speech, and Signal Processing}, vol. 32,
no. 2, pp. 236--243, 1984.

\bibitem{mulan}
Q. Huang et al.,
\textit{MuLan: A Joint Embedding of Music Audio and Natural Language},
Proceedings of the International Society for Music Information Retrieval
Conference, 2022.

\bibitem{muq}
H. Zhu, Y. Zhou, H. Chen, J. Yu, Z. Ma, R. Gu, et al.,
``MuQ: Self-Supervised Music Representation Learning with Mel Residual Vector
Quantization,''
\textit{IEEE Transactions on Audio, Speech and Language Processing}, 2025.

\bibitem{ddpm}
J. Ho, A. Jain, and P. Abbeel,
``Denoising Diffusion Probabilistic Models,''
\textit{Advances in Neural Information Processing Systems}, vol. 33,
pp. 6840--6851, 2020.

\bibitem{flowmatching}
Y. Lipman, R. T. Q. Chen, H. Ben-Hamu, M. Nickel, and M. Le,
``Flow Matching for Generative Modeling,''
\textit{International Conference on Learning Representations}, 2023.

\bibitem{rectifiedflow}
X. Liu, C. Gong, and Q. Liu,
``Flow Straight and Fast: Learning to Generate and Transfer Data with Rectified
Flow,''
\textit{International Conference on Learning Representations}, 2023.

\bibitem{dit}
W. Peebles and S. Xie,
``Scalable Diffusion Models with Transformers,''
\textit{Proceedings of the IEEE/CVF International Conference on Computer Vision},
pp. 4195--4205, 2023.

\bibitem{distillation}
G. Hinton, O. Vinyals, and J. Dean,
\textit{Distilling the Knowledge in a Neural Network},
arXiv:1503.02531, 2015.

\bibitem{vocos}
H. Siuzdak,
``Vocos: Closing the Gap between Time-Domain and Fourier-Based Neural Vocoders
for High-Quality Audio Synthesis,''
\textit{International Conference on Learning Representations}, vol. 2024,
pp. 25719--25733, 2024.

\bibitem{bigvgan}
S.-g. Lee, W. Ping, B. Ginsburg, B. Catanzaro, and S. Yoon,
``BigVGAN: A Universal Neural Vocoder with Large-Scale Training,''
\textit{International Conference on Learning Representations}, 2023.

\bibitem{dpo}
R. Rafailov, A. Sharma, E. Mitchell, S. Ermon, C. D. Manning, and C. Finn,
``Direct Preference Optimization: Your Language Model Is Secretly a Reward
Model,''
\textit{Advances in Neural Information Processing Systems}, vol. 36, 2023.

\bibitem{xlmr}
A. Conneau et al.,
``Unsupervised Cross-lingual Representation Learning at Scale,''
\textit{Proceedings of ACL}, pp. 8440--8451, 2020.

\bibitem{xphonebert}
L. T. Nguyen, T. Pham, and D. Q. Nguyen,
``XPhoneBERT: A Pre-trained Multilingual Model for Phoneme Representations for
Text-to-Speech,''
arXiv preprint arXiv:2305.19709, 2023.

\bibitem{whisper}
A. Radford et al.,
``Robust Speech Recognition via Large-Scale Weak Supervision,''
\textit{Proceedings of ICML}, pp. 28492--28518, 2023.

\bibitem{demucs}
A. D\'efossez, N. Usunier, L. Bottou, and F. Bach,
``Music Source Separation in the Waveform Domain,''
arXiv:1911.13254, 2019.

\bibitem{htdemucs}
S. Rouard, F. Massa, and A. D\'efossez,
``Hybrid Transformers for Music Source Separation,''
in \textit{ICASSP 2023--2023 IEEE International Conference on Acoustics, Speech
and Signal Processing (ICASSP)}, pp. 1--5, IEEE, 2023.


\bibitem{diffrhythm2}
Y. Jiang, H. Chen, Z. Ning, J. Yao, Z. Han, D. Wu, et al.,
``DiffRhythm 2: Efficient and High Fidelity Song Generation via Block Flow
Matching,''
arXiv preprint arXiv:2510.22950, 2025.


\end{thebibliography}

\end{document}
